% SPDX-License-Identifier: Apache-2.0
%
% Datasheets: one per component the project has built. Built by
% //docs:datasheets. The tables of ports, state and children come from
% each component's own lowering, through //tools/datasheet; the prose is
% in docs/datasheets/, one file per component, and
% //tools/datasheet:coverage_test fails when a component has none.
\documentclass[journal,onecolumn,9pt]{IEEEtran}

\newcommand{\listingsize}{\normalsize}
% SPDX-License-Identifier: Apache-2.0
%
% The house preamble, shared by every document in this directory. Loaded
% after \documentclass, before \title. Keeping it in one file is what
% stops two articles drifting apart in font, listing style or figure
% style.
% Computer Modern. IEEEtran selects Times (ptm) by default, so the face has
% to be chosen explicitly.
%
% Latin Modern is the Computer Modern variety used here, and the choice is
% forced rather than aesthetic. Under T1 encoding, plain `cmr` has no Type 1
% outlines unless cm-super is installed, and it is not in the pinned
% distribution, so pdflatex falls back to EC bitmaps and every glyph in the
% PDF becomes a Type 3 font. That was measured: the first build produced a
% document with 10 Type 3 fonts and no Type 1. Latin Modern is Computer
% Modern redrawn with a full T1 repertoire and real .pfb outlines, so the
% same design comes out scalable.
\usepackage[T1]{fontenc}
\usepackage[utf8]{inputenc}
\usepackage{lmodern}
% microtype is not in the pinned TeX distribution. emergencystretch alone
% keeps the two-column measure from producing overfull lines.
\setlength{\emergencystretch}{3em}

\usepackage{amsmath}
\usepackage{amssymb}
\usepackage{amsthm}
\usepackage{booktabs}
\usepackage{array}
% enumitem is not in the pinned TeX distribution; the list spacing below
% does what \setlist would have done.
\usepackage{float}
\usepackage{xcolor}
\usepackage{listings}
\usepackage{tikz}
\usetikzlibrary{arrows.meta,positioning,fit,backgrounds,calc}
\usepackage[hidelinks,bookmarks=true]{hyperref}

% Numbered, definition-style box for a rule the reader is meant to apply.
\theoremstyle{definition}
\newtheorem{rulebox}{Rule}
\newtheorem{example}{Example}

% Unnumbered asides.
\theoremstyle{remark}
\newtheorem*{tip}{Tip}
\newtheorem*{pitfall}{Pitfall}

% Tighter lists, in place of enumitem's \setlist.
\setlength{\topsep}{2pt}
\setlength{\itemsep}{1pt}
\setlength{\parsep}{0pt}

\newcommand{\code}[1]{\texttt{#1}}
\newcommand{\term}[1]{\emph{#1}}

% Syntax highlighting. Four roles, distinguishable in colour and still
% distinguishable in greyscale, because an IEEE article gets printed: the
% keyword is the only bold role, the comment is the only italic one, and the
% two remaining roles differ in lightness as well as in hue.
\definecolor{kwblue}{RGB}{20,60,150}
\definecolor{cmtgreen}{RGB}{90,120,90}
\definecolor{strred}{RGB}{150,50,40}
\definecolor{numgray}{gray}{0.45}
\definecolor{framegray}{gray}{0.70}
\definecolor{bgshade}{gray}{0.97}
% A darker fill for figure cells. The listing background has to stay very
% light to keep code readable; a figure cell has to be clearly shaded or the
% distinction it draws is invisible in print.
\definecolor{cellshade}{gray}{0.90}

% The listing size is the document's choice, because it depends on the
% column: a two-column article sets `\listingsize` to 6pt and keeps its
% listings within 70 columns, a one-column document keeps the body size
% and includes source files at 80. Lines are never broken; a line that does not fit
% overflows the frame, which is visible, rather than wrapping, which is
% not.
\providecommand{\listingsize}{\scriptsize}

% `'` is deliberately not a string delimiter: the language uses it for loop
% labels, as Rust does, so treating it as a quote would swallow the rest of
% the line.
\lstdefinestyle{txhdl}{
  basicstyle=\ttfamily\listingsize,
  keywordstyle=\bfseries\color{kwblue},
  commentstyle=\itshape\color{cmtgreen},
  stringstyle=\color{strred},
  numberstyle=\tiny\color{numgray},
  morekeywords={mod,use,pub,const,type,struct,enum,trait,impl,fn,async,let,mut,
    match,if,else,loop,while,for,in,break,continue,return,as,await,
    module,bus,channel,transaction,protocol,pipeline,flow,seq,config,
    port,stage,pipe,instance,connect,bind,tag,domain,blackbox,foreign,
    property,role,signal,handler,true,false,
    sequence,interface,modport,implements,package,import,var,wire,func,proc,
    case,default,next,fork,branch,join,state,fsm},
  morecomment=[l]{//},
  morestring=[b]",
  showstringspaces=false,
  columns=fullflexible,
  keepspaces=true,
  breaklines=false,
  % Framed, so a listing is bounded rather than floating in the column.
  frame=single,
  rulecolor=\color{framegray},
  framesep=3pt,
  backgroundcolor=\color{bgshade},
  xleftmargin=3pt,
  xrightmargin=3pt,
  aboveskip=6pt,
  belowskip=6pt,
  % Every listing is numbered and captioned, so it can be referred to by
  % number. `caption` without `float` numbers the listing and leaves it
  % where it was written, which is what a listing in running prose needs.
  captionpos=b,
  abovecaptionskip=2pt,
  belowcaptionskip=6pt,
}

% Figures drawn as text. Framed the same way, and with the highlighting
% switched off, because a box-and-arrow diagram is not code and half its
% words turning blue reads as an accident. Setting a style adds keys rather
% than clearing the ones already set, so the three roles are neutralised by
% hand rather than by leaving them out.
\lstdefinestyle{ascii}{
  basicstyle=\ttfamily\listingsize,
  keywordstyle=\ttfamily,
  commentstyle=\ttfamily,
  stringstyle=\ttfamily,
  columns=fullflexible,
  keepspaces=true,
  frame=single,
  rulecolor=\color{framegray},
  framesep=3pt,
  backgroundcolor=\color{bgshade},
  xleftmargin=3pt,
  xrightmargin=3pt,
  aboveskip=6pt,
  belowskip=6pt,
}
\lstset{style=txhdl}

% House TikZ styles. `shorten` lives on the arrow styles rather than on each
% arrow, so every arrow in the document gets the gap, including ones added
% later. TikZ clips a path at a node's border, so without it an arrowhead
% butts against the box with no gap at all. Keep `node distance` well above
% twice the shorten, or the arrow shrinks to nothing.
\tikzset{
  box/.style={draw,rounded corners=2pt,align=center,inner sep=4pt,
              font=\footnotesize},
  boxf/.style={box,fill=bgshade},
  lbl/.style={font=\scriptsize\itshape,align=center},
  fl/.style={-{Stealth[length=4pt]},shorten >=2.5pt,shorten <=2.5pt},
  fld/.style={fl,dashed},
}



\InputIfFileExists{buildstamp.tex}{}{}
\providecommand{\buildstamp}{Draft build (outside the Bazel flow).}

% The generated tables. A sheet that asks for a table the generator
% does not write stops the build, rather than printing nothing.
\input{tools/datasheet/ds_tables}
\makeatletter
\newcommand{\dsget}[2]{%
  \@ifundefined{ds@#1@#2}%
    {\PackageError{datasheets}{//tools/datasheet writes no #1 for #2}%
      {Add the component to tools/datasheet/main.rs.}}%
    {\csname ds@#1@#2\endcsname}}
\makeatother

% A datasheet's heading: the component and what it is, on a page of
% its own.
\newcommand{\datasheet}[2]{\clearpage\section{#1}\label{ds:#1}%
\noindent\textit{#2}\par\medskip}

% Where a component lives: its source, its Rust path, and the document
% that explains it.
\newcommand{\dsfacts}[3]{%
\begin{center}\footnotesize
\begin{tabular}{@{}ll@{}}
Source & #1 \\
Rust path & #2 \\
Explained in & #3 \\
\end{tabular}
\end{center}}

% The generated part of a sheet.
\newcommand{\dstables}[1]{%
\subsection*{Ports, state and children}
As lowered for \dsget{type}{#1}: \dsget{counts}{#1}.
\dsget{ports}{#1}
\dsget{state}{#1}
\dsget{children}{#1}}

\title{TxHDL Datasheets}

\author{Filip Filmar and Dragiša Janković%
\thanks{Both authors are independent. Filip Filmar is the
corresponding author, at f@hdlfactory.com; Dragiša Janković is
reached at linkedin.com/in/dragisa-jankovic.}%
\thanks{One datasheet per component built in the project repository
\code{HDL/txhdl}. The tables of ports, state and children on each sheet
are generated at build time from the component's own lowering, so they
are the netlist's. The cover, \code{//docs:cover}, lists every
document.}%
\thanks{The authors used a large language model, Claude, as an
assistant in exploring the concepts and in writing and constructing
the documents and the programs; every commit of the repository says
so and carries the prompts that led to it, verbatim.}%
\thanks{\buildstamp}}

\markboth{TxHDL Datasheets}{Filmar and Janković: TxHDL Datasheets}

\begin{document}
\maketitle

\section*{How to Read a Datasheet}

A component is a unit written in the lowered subset, so it is both a
simulation and a netlist. Each sheet says what the component does, the
parameters it takes, how it behaves at its ports, how it is checked and
where it is used, and what it does not do.

The tables of ports, state and children are not typed in. The build
lowers each component with the parameters the tree uses it with, the
ones the sheet names, and writes the tables from that lowering. A
channel port is one row, and its width is the width of the word it passes;
in the netlist it is \code{\_data}, \code{\_valid} and \code{\_ready}. A
wire port is one row and one net. A memory's width is its word's, and
its words are how many it holds.

A family of components written by one macro, such as the routers of two
to eight peripherals, has one sheet, and its tables are those of the
member the tree uses.

% Channels and queues.
% SPDX-License-Identifier: Apache-2.0
% covers: Buffer
\datasheet{Buffer}{The channel as hardware: an elastic buffer of two.}

\dsfacts{\code{lib/parts/src/buffer.rs}}{\code{txhdl\_parts::buffer::Buffer}}%
{\code{//docs:runtime}, \code{//docs:examples}}

\subsection*{Function}

A channel in a simulation is a buffer of two between a sender and a
receiver. \code{Buffer} is that buffer as a netlist, for the places
where two lowered units meet on a channel: a board top, or a design
wired by hand. Its sender side is \code{tx\_data}, \code{tx\_valid} and
\code{tx\_ready}; its receiver side is \code{rx\_data}, \code{rx\_valid}
and \code{rx\_ready}.

\subsection*{Parameters}

\begin{center}\footnotesize
\begin{tabular}{@{}lp{0.7\textwidth}@{}}
\toprule
Parameter & Meaning \\
\midrule
\code{W} & The width of the words the channel holds, in bits. \\
\bottomrule
\end{tabular}
\end{center}

\dstables{Buffer}

\subsection*{Behaviour}

\begin{itemize}
\item \code{tx\_ready} is high when the tail is empty, as the last edge
left it. \code{rx\_valid} is high when the head is full, and
\code{rx\_data} is the head.
\item All three outputs are registers, so there is no combinational
path from one side to the other.
\item At an edge, a take moves the tail into the head, and a push goes
into the head if the head is then empty, else into the tail.
\item With both sides running every cycle, one word passes per cycle.
\end{itemize}

\subsection*{Verification}

\begin{itemize}
\item \code{buffer\_is\_the\_runtime\_channel} in \code{txhdl\_parts}
runs the buffer beside a runtime channel for 200 cycles of pseudorandom
offers and takes, and checks \code{ready}, \code{valid} and the head
after every cycle.
\item \code{ex\_buffer} does the same on a fixed pattern, and the build
simulates \code{Buffer<8>}'s netlist against that run under nvc and
Verilator: \code{//docs:buffer\_sim\_buffer8\_tb\_test} and
\code{//docs:buffer\_vsim\_test}.
\end{itemize}

\subsection*{Used by}

The example \code{ex\_buffer}. A unit of units does not need it, since
its netlist puts a \code{txhdl\_chan} of the same rules between its
children.

\subsection*{Limits}

It holds two words. A consumer that takes in bursts needs \code{Fifo}.

% SPDX-License-Identifier: Apache-2.0
% covers: Fifo
\datasheet{Fifo}{A FIFO of a power of two words, with a channel at each end.}

\dsfacts{\code{lib/parts/src/fifo.rs}}{\code{txhdl\_parts::fifo::Fifo}}%
{\code{//docs:runtime}, \code{//docs:examples}, \code{//docs:station}}

\subsection*{Function}

A channel is a buffer of two. \code{Fifo} is the depth a channel does
not have, for a unit whose consumer takes in bursts or holds off. It
has a channel in, \code{inp}, a channel out, \code{out}, and \code{D}
words between. \code{//docs:station} puts one behind a reservation
station for that reason.

\subsection*{Parameters}

\begin{center}\footnotesize
\begin{tabular}{@{}lp{0.7\textwidth}@{}}
\toprule
Parameter & Meaning \\
\midrule
\code{T} & What it holds: a \code{Transaction} and \code{Value} type. \\
\code{AW} & The width of its pointers, in bits. \\
\code{D} & How many words it holds, which must be \code{1 << AW}. \\
\bottomrule
\end{tabular}
\end{center}

\code{D} is stated rather than computed, because an expression in a
const parameter's position needs nightly Rust. The tables below use
\code{Fifo<U<8>, 2, 4>}, the FIFO of \code{ex\_queue}: four bytes.

\dstables{Fifo}

\subsection*{Behaviour}

\begin{itemize}
\item \code{mem} holds the words. \code{head} is where the oldest word
is, and \code{tail} is where the next word lands.
\item The pointers are equal both when it is empty and when it is
full. \code{full} says which.
\item A word is taken from \code{inp} whenever \code{full} is clear,
and written at \code{tail}, which moves on by one.
\item When it is not empty and \code{out} has room, the word at
\code{head} is sent on \code{out}, and \code{head} moves on by one.
\item A word taken lands in the memory at the end of its step and is
offered the step after. The source states the latency as one cycle.
\item \code{full} clears in a cycle that sends a word. It sets in a
cycle that takes a word and sends none, when the moved \code{tail}
meets \code{head}.
\end{itemize}

\subsection*{Verification}

\begin{itemize}
\item \code{fifo\_keeps\_every\_word\_in\_order} in
\code{txhdl\_parts}, in \code{//lib/parts:parts\_test}, runs
\code{Fifo<U<8>, 2, 4>} against a queue. It offers words at random
for 200 cycles and takes none for the first 40, so the FIFO fills. It
checks every word comes out in order, that \code{full} was seen, and
that it is empty at the end.
\item \code{ex\_queue} bursts nine words into the same FIFO while the
sink takes nothing for eight cycles, as \code{//docs:examples}
describes. The build simulates the netlist, named \code{queue4},
against that run under nvc and Verilator:
\code{//docs:queue\_sim\_queue4\_tb\_test} and
\code{//docs:queue\_vsim\_test}.
\end{itemize}

\subsection*{Used by}

The example \code{ex\_queue}. \code{ex\_fifo} is a different unit: an
asynchronous FIFO of its own, defined in the example.

\subsection*{Limits}

It runs on one clock, \code{DefaultClock}. \code{D} must be
\code{1 << AW}, and nothing checks it. When \code{full} is set it
takes no word, even in a cycle in which it sends one.

% SPDX-License-Identifier: Apache-2.0
% covers: Station2, Station3, Station4, Station5, Station6, Station7, Station8, Station9, Station10
\datasheet{Station}{Reservation stations of two to ten inputs: operands gathered by tag, sent as one line.}

\dsfacts{\code{lib/parts/src/station.rs}; generator \code{station\_text} in \code{lib/macros/lib.rs}}%
{\code{txhdl\_parts::station::Station2} to \code{Station10}}%
{\code{//docs:station}, \code{//docs:runtime}, \code{//docs:examples}}

\subsection*{Function}

A reservation station gathers the operands of an operation that arrive
on separate channels, out of order, and fires when all of them have.
Each input is a channel of \code{Tagged<TB, T>}: a tag of \code{TB}
bits above a value. The station has a line per tag value, with a cell
per input. An input lands in its cell of the line its tag names.

A line whose every cell holds is sent on the one output as a
\code{Line}\textit{N}: the tag, and a value per input, \code{v0}
upwards. \code{station!(StationN,
N)} writes each station with its line struct beside it, since the
lowering reads a body and not a loop over inputs. \code{station.rs}
writes \code{Station2} to \code{Station10}; the macro itself accepts 2
to 16.

\subsection*{Parameters}

\begin{center}\footnotesize
\begin{tabular}{@{}lp{0.7\textwidth}@{}}
\toprule
Parameter & Meaning \\
\midrule
\code{TB} & The tag width, in bits. It is two at least. \\
\code{L} & The count of lines, which is \code{1 << TB}. \\
\code{T}\textit{i} & The value type of input \textit{i}, for \textit{i}
from 0 to \textit{N} less one. The types may differ. \\
\bottomrule
\end{tabular}
\end{center}

\code{L} is stated rather than computed, because a computed width in a
type needs nightly Rust. The tag has two bits at least, since a one-bit
value is a single wire in VHDL and cannot index a line. The tables
below use \code{Station3<4, 16, U<2>, U<16>, U<16>>}, the station of
\code{ex\_compose}: sixteen lines, an operation and two operands.

\dstables{Station}

\subsection*{Behaviour}

Each input \textit{i} has \code{mem}\textit{i}, its cells, a word per
line, and \code{occ}\textit{i}, a bit per line that says which cells
hold. Each cycle:

\begin{itemize}
\item An input's cell is free when its tag's bit of the input's mask
is clear.
\item Taking an input completes its line when, for every other input,
that line's cell holds, or that input offers the same tag now and its
own cell is free.
\item The line sent is the completing input's of lowest index, if any
input completes and the output has room. At most one line leaves per
cycle.
\item An input is taken when its cell is free and taking it completes
nothing, or completes the line being sent. An input's \code{ready} is
that take.
\item So an input never waits on another input's cell. It waits on its
own cell, on the output's room, or while an input of lower index
completes a different line in the same cycle.
\item A taken input writes its cell and sets its bit. The line sent
clears its bit in every mask.
\item The values sent come from the cells. The value of an input that
completes the line in this cycle goes straight through instead.
\end{itemize}

\subsection*{Verification}

\begin{itemize}
\item Two tests in \code{txhdl\_parts}, in
\code{//lib/parts:parts\_test}. \code{station2\_follows\_the\_rule} runs
\code{Station2<2, 4, U<8>, U<16>>} against the rule written with loops
for 400 cycles of random offers and takes. It checks both masks after
every cycle and every line the sink takes.
\code{station3\_lowest\_input\_wins} sets up two lines that complete in
one cycle, and checks that the lower input's line goes first.
\item \code{ex\_station} runs a \code{Station2<2, 4, U<8>, U<16>>}
beside \code{station.v}, the same station written by hand in Verilog,
and compares every line the two send. The build simulates the
station's netlist against that run under nvc and Verilator:
\code{//docs:station\_sim\_station\_tb\_test} and
\code{//docs:station\_vsim\_test}.
\item \code{ex\_compose} lowers \code{Top}, which holds the
\code{Station3} of the tables and the stream ALU as children. The build
simulates the netlist of the top against that run:
\code{//docs:compose\_sim\_top\_tb\_test} and
\code{//docs:compose\_vsim\_test}.
\end{itemize}

\subsection*{Used by}

\code{Station2} in the example \code{ex\_station}. \code{Station3} in
the example \code{ex\_compose}, in front of the stream ALU, whose
requests are \code{Line3} values.

\subsection*{Limits}

A station is not an arbiter. If an input's next offer names a cell it
already holds, it waits until that line completes. If the other inputs
never offer that tag, it waits for ever, as \code{//docs:station}
states. Each input must, in time, see every tag its lines wait on. One
line leaves per cycle at most, and a line waits for room on the output;
\code{//docs:station} puts a \code{Fifo} behind the station where the
consumer takes in bursts.

% The AXI link.
% SPDX-License-Identifier: Apache-2.0
% covers: AxiHost
\datasheet{AxiHost}{The host side of an AXI4 link: bursts out, identifiers handed out, answers back.}

\dsfacts{\code{lib/parts/src/bus/axi.rs}}{\code{txhdl\_parts::bus::axi::AxiHost}}%
{\code{//docs:axi}, \code{//docs:noc}}

\subsection*{Function}

An AXI4 link in \code{txhdl\_parts} has a tracker at each end, between
a client and the five AXI channels. \code{AxiHost} is the tracker on
the host side. A host client gives it an \code{Issue} on \code{issue}:
an address phase with no identifier, and a bit that says whether it is
a read. \code{AxiHost} gives the burst an identifier, sends it on
\code{ar} or \code{aw}, and tells the client the identifier on
\code{grant}.

It also passes the client's write beats from \code{wbeat} to \code{w},
turns each \code{B} beat into a \code{Done} on \code{done}, and passes
each read beat from \code{r} to \code{rdata}. The client hands an
identifier back on \code{release} when it has taken that burst's
answer. \code{axi}, \code{axi\_units} and \code{axi\_to\_unit} make the
channels it runs on. Its ports are named types, \code{HostIn} and
\code{HostOut}, so a design may supply a unit of its own instead.

\subsection*{Parameters}

\begin{center}\footnotesize
\begin{tabular}{@{}lp{0.7\textwidth}@{}}
\toprule
Parameter & Meaning \\
\midrule
\code{A} & The address width, in bits. \\
\code{D} & The data width, in bits. \\
\code{S} & The write strobe width, one bit per byte lane, so \code{D / 8}. \\
\code{I} & The identifier width, in bits. \\
\code{NIDS} & How many identifiers it hands out, so \code{1 << I}. It is
the width of \code{busy}, and the most bursts in flight at once. \\
\bottomrule
\end{tabular}
\end{center}

\code{S} and \code{NIDS} are stated rather than computed, because an
expression in a const parameter's position needs nightly Rust. The
source states that a link whose \code{S} and \code{NIDS} disagree with
its \code{D} and \code{I} will not behave. The tables below use
\code{AxiHost<32, 32, 4, 2, 4>}, the tracker of Vreteno's link and of
its board.

\dstables{AxiHost}

\subsection*{Behaviour}

\begin{itemize}
\item At each rising edge it looks at the issue at the head of
\code{issue}. The identifier that burst would take is \code{turn}.
\item The burst goes out when that identifier's bit of \code{busy} is
clear, \code{grant} has room, and the address channel it needs has
room: \code{ar} for a read, \code{aw} for a write.
\item A burst that goes out is sent with \code{turn} as its \code{id}.
Its \code{addr}, \code{len}, \code{size}, \code{burst}, \code{lock},
\code{cache}, \code{prot}, \code{qos} and \code{region} are the
issue's.
\item In that cycle \code{grant} sends the identifier, \code{turn}
moves on by one, and the identifier's bit of \code{busy} is set. At
most one burst goes out per cycle.
\item The identifiers go round robin. A burst waits for its own turn's
identifier to be free, and for nothing else. Another free identifier
does not let it go.
\item \code{release} is taken whenever it offers, and it clears the
bit of the identifier it names. An identifier is therefore outstanding
from its address phase until the client has taken its answer.
\code{//docs:axi} explains why freeing it on the response would be
wrong.
\item A write beat passes from \code{wbeat} to \code{w} unchanged when
\code{w} has room. AXI4 puts no identifier on \code{w}, so the source
says a client sends a burst's beats before it issues the next burst.
\item A \code{B} beat becomes a \code{Done} of the same \code{id} and
\code{resp} when \code{done} has room. A read beat passes from
\code{r} to \code{rdata} unchanged when \code{rdata} has room.
\end{itemize}

\subsection*{Verification}

\begin{itemize}
\item Four tests in \code{bus::axi} run \code{AxiHost<16, 32, 4, 2, 4>}
and \code{AxiPer<16, 32, 4, 2>} with a host client and a peripheral
client on the two ends:
\begin{itemize}
\item \code{a\_read\_reads\_what\_the\_write\_wrote};
\item \code{answers\_come\_back\_out\_of\_order};
\item \code{a\_transaction\_is\_answered\_by\_another\_process};
\item \code{every\_burst\_is\_answered\_and\_reads\_agree\_with\_writes},
which issues 24 bursts of one to four words to a peripheral that
\code{serve} runs with four slots.
\end{itemize}
\item The tests of \code{bus::axi::sim} put the same two trackers in
front of the simulated RAM:
\begin{itemize}
\item \code{a\_burst\_reads\_what\_a\_burst\_wrote};
\item \code{a\_strobe\_covers\_the\_lanes\_it\_says};
\item \code{a\_fixed\_burst\_stays\_on\_one\_word};
\item \code{an\_address\_past\_the\_end\_is\_answered};
\item \code{an\_image\_is\_loaded\_and\_read\_back}.
\end{itemize}
All of these run in \code{//lib/parts:parts\_test}.
\item \code{ex\_axi} issues three bursts and has them answered out of
order. The build simulates the netlist of \code{AxiHost<16, 32, 4, 2,
4>} against that run under nvc and Verilator:
\begin{itemize}
\item \code{//docs:axi\_sim\_axi\_host\_axi\_host\_tb\_test};
\item \code{//docs:axi\_vsim\_axi\_host\_test}.
\end{itemize}
\item The tests of the router, the AXI-Lite bridge, the Wishbone bridge
and the network on chip in \code{txhdl\_parts} run it in simulation.
So do \code{//ddr3:ddr3\_test}, \code{//gpu/razboj:razboj\_test},
\code{//soc:demo\_test} and \code{//cpu/vreteno:lockstep}.
\end{itemize}

\subsection*{Used by}

The examples \code{ex\_axi}, \code{ex\_axi\_serve}, \code{ex\_axi\_router},
\code{ex\_axi\_lite}, \code{ex\_axi\_reg}, \code{ex\_axi\_wb},
\code{ex\_eth} and \code{ex\_hdmi}. Vreteno's run in
\code{cpu/vreteno/src/run.rs} and \code{cpu/vreteno/src/main.rs}, and
the board's design, \code{Board}, as its field \code{host}. The
four-corner system in \code{soc/src/lib.rs}, once for the core and
once for Razboj. Razboj's simulation in \code{gpu/razboj/src/sim.rs}.

\subsection*{Limits}

At most \code{NIDS} bursts are in flight at once. A client that never
releases an identifier stops the unit when the turn reaches that
identifier. It sends \code{lock}, \code{cache}, \code{qos} and
\code{region} as the client gave them and acts on none of them; nothing
in \code{bus::axi} makes an exclusive access. It does not check
\code{S} and \code{NIDS} against \code{D} and \code{I}.

% SPDX-License-Identifier: Apache-2.0
% covers: AxiPer
\datasheet{AxiPer}{The peripheral side of an AXI4 link: two address channels merged into one stream of requests.}

\dsfacts{\code{lib/parts/src/bus/axi.rs}}{\code{txhdl\_parts::bus::axi::AxiPer}}%
{\code{//docs:axi}, \code{//docs:vreteno}, \code{//docs:noc}}

\subsection*{Function}

\code{AxiPer} is the tracker at the peripheral end of an AXI4 link. It
merges \code{aw} and \code{ar} into one stream of \code{PerReq} on
\code{req}, each with a bit that says whether it is a read. It passes
the write beats from \code{w} to \code{wd}. It turns the client's
\code{Answer} on \code{ans} into \code{B} beats on \code{b}, and
passes the client's read beats from \code{rb} to \code{r}.

It allocates no identifiers. A response leaves with the identifier the
client gave it. Its ports are named types, \code{PerIn} and
\code{PerOut}, so a design may supply a unit of its own instead.
\code{//docs:axi} notes that a peripheral may also read the five AXI
channels directly and leave \code{AxiPer} out.

\subsection*{Parameters}

\begin{center}\footnotesize
\begin{tabular}{@{}lp{0.7\textwidth}@{}}
\toprule
Parameter & Meaning \\
\midrule
\code{A} & The address width, in bits. \\
\code{D} & The data width, in bits. \\
\code{S} & The write strobe width, one bit per byte lane, so \code{D / 8}. \\
\code{I} & The identifier width, in bits. \\
\bottomrule
\end{tabular}
\end{center}

There is no \code{NIDS}, since this end hands out no identifiers. The
tables below use \code{AxiPer<32, 32, 4, 2>}, the tracker in front of
Vreteno's data memory and timer.

\dstables{AxiPer}

\subsection*{Behaviour}

\begin{itemize}
\item At most one request leaves per cycle, and only when \code{req}
has room.
\item \code{rr} says which address channel goes first: zero for
\code{ar}. When both offer, the one that goes first is taken. When
only one offers, that one is taken.
\item Taking a read sets \code{rr} to one, and taking a write sets it
to zero. So two channels that both stream take turns, and neither
starves the other.
\item The request has \code{read} high when it came from \code{ar}.
Its \code{id}, \code{addr}, \code{len}, \code{size}, \code{burst},
\code{lock}, \code{cache}, \code{prot}, \code{qos} and \code{region}
are the address phase's.
\item A write beat passes from \code{w} to \code{wd} unchanged when
\code{wd} has room, independently of the requests.
\item An answer on \code{ans} becomes a \code{B} of the same \code{id}
and \code{resp} when \code{b} has room. A read beat passes from
\code{rb} to \code{r} unchanged when \code{r} has room.
\end{itemize}

\subsection*{Verification}

\begin{itemize}
\item Four tests in \code{bus::axi} run \code{AxiHost<16, 32, 4, 2, 4>}
and \code{AxiPer<16, 32, 4, 2>} with a client on each end:
\begin{itemize}
\item \code{a\_read\_reads\_what\_the\_write\_wrote};
\item \code{answers\_come\_back\_out\_of\_order};
\item \code{a\_transaction\_is\_answered\_by\_another\_process};
\item \code{every\_burst\_is\_answered\_and\_reads\_agree\_with\_writes}.
\end{itemize}
\item \code{bus::axi::fuzz} runs the two trackers with a relay that
stalls at random in every channel and a seeded client and peripheral
on the ends, and checks every burst answered once, reads against
writes, \code{last} on the final beat, and no deadlock, on 24 seeds;
\code{//lib/parts:axi\_fuzz\_test}, manual, runs more.
\item The tests of \code{bus::axi::sim} put the two trackers in front
of the simulated RAM:
\begin{itemize}
\item \code{a\_burst\_reads\_what\_a\_burst\_wrote};
\item \code{a\_strobe\_covers\_the\_lanes\_it\_says};
\item \code{a\_fixed\_burst\_stays\_on\_one\_word};
\item \code{an\_address\_past\_the\_end\_is\_answered};
\item \code{an\_image\_is\_loaded\_and\_read\_back}.
\end{itemize}
All of these run in \code{//lib/parts:parts\_test}.
\item \code{ex\_axi} has three bursts answered out of order. The build
simulates the netlist of \code{AxiPer<16, 32, 4, 2>} against that run
under nvc and Verilator:
\begin{itemize}
\item \code{//docs:axi\_sim\_axi\_per\_axi\_per\_tb\_test};
\item \code{//docs:axi\_vsim\_axi\_per\_test}.
\end{itemize}
\item The tests of the router, the Wishbone bridge and the network on
chip in \code{txhdl\_parts} run it in simulation. So do
\code{//ddr3:ddr3\_test}, \code{//gpu/razboj:razboj\_test},
\code{//soc:demo\_test} and \code{//cpu/vreteno:lockstep\_test}.
\end{itemize}

\subsection*{Used by}

The examples \code{ex\_axi}, \code{ex\_axi\_serve},
\code{ex\_axi\_router}, \code{ex\_axi\_reg} and \code{ex\_axi\_wb}.
Vreteno's run in \code{cpu/vreteno/src/run.rs} and
\code{cpu/vreteno/src/main.rs}, in front of the data memory and the
timer. The board's design, \code{Board}, as its fields \code{pdmem},
\code{ptimer} and \code{pddr3}. The four-corner system in
\code{soc/src/lib.rs}, and Razboj's simulation in
\code{gpu/razboj/src/sim.rs}.

\subsection*{Limits}

It does not count beats, look at \code{last}, or match write beats to
requests. A client does that: \code{Per::accept} gathers a write's
beats after its request. It does not check that an answer's identifier
belongs to an open request. It passes \code{lock}, \code{cache},
\code{qos} and \code{region} to the client and acts on none of them.

% SPDX-License-Identifier: Apache-2.0
% covers: Router
\datasheet{Router}{The AXI4 router: one host and N peripherals, decoded by an address map.}

\dsfacts{\code{lib/parts/src/bus/router.rs}}%
{\code{txhdl\_parts::bus::router::Router}}%
{\code{//docs:axi}, \code{//docs:vreteno}, \code{//docs:noc}}

\subsection*{Function}

A router puts several peripherals behind one host. It sits in the five
channels of an AXI4 link, so nothing on either side changes when it is
put there. It takes the host's \code{aw}, \code{ar} and \code{w}, and
gives each peripheral \textit{i} its own, element \textit{i} of the
arrays \code{aws}, \code{ars} and \code{ws}. It takes each peripheral's
answers from \code{bs} and \code{rs}, and gives the host one \code{b}
and one \code{r}. In the netlist element \textit{i} of an array port is
a port of its own, \code{aws\_}\textit{i} and so on.

One unit serves every count of peripherals: the peripheral sides are
arrays of \code{N}, and the lowering unrolls the loops over them when
\code{lowered} runs (issue 500). Each peripheral's range is its entry
of the map, a type \code{M} implementing \code{txhdl::map::AddrMap<N>}
whose constant \code{RANGES} is \code{N} pairs of a base and a mask,
read the same way (issue 593). A burst whose address is in no range is
answered \code{DecErr} by the router itself. The router was a family of
units, \code{Router2} to \code{Router8}, written by \code{router!},
until \#648.

\subsection*{Parameters}

\begin{center}\footnotesize
\begin{tabular}{@{}lp{0.7\textwidth}@{}}
\toprule
Parameter & Meaning \\
\midrule
\code{N} & The number of peripherals. \\
\code{M} & The address map: a type implementing \code{AddrMap<N>},
whose \code{RANGES[i]} is peripheral \textit{i}'s base and mask. It
holds no signal; the router keeps it as \code{PhantomData}. \\
\code{A} & The address width, in bits. \\
\code{D} & The data width, in bits. \\
\code{S} & The write strobe width, one bit per byte lane, so \code{D / 8}. \\
\code{I} & The identifier width, in bits. \\
\bottomrule
\end{tabular}
\end{center}

\subsection*{Address map}

An address \code{a} is peripheral \textit{i}'s when
\code{(a \& mask) == base} for \code{RANGES[i]} and for no range before
it: the first range that matches decides. The router decodes the
address of each address phase. Every address in no range is a hole.
The tables below are generated for the board's router,
\code{BoardRouter} in \code{cpu/vreteno/src/board.rs}, which is
\code{Router<7, BoardMap, 32, 32, 4, 4>} with \code{BoardMap} beside
it:

\begin{center}\footnotesize
\begin{tabular}{@{}llll@{}}
\toprule
Peripheral & Base & Mask & On the board \\
\midrule
0 & \code{0x1000} & \code{0xffff\_f000} & The data memory \\
1 & \code{0x0200\_0000} & \code{0xffff\_0000} & The timer \\
2 & \code{0x3000} & \code{0xffff\_f000} & The AXI-Lite bridge to the
serial port and its neighbours \\
3 & \code{0x4000\_0000} & \code{0xc000\_0000} & The DDR3 memory \\
4 & \code{0x0c00\_0000} & \code{0xfc00\_0000} & The interrupt controller \\
5 & \code{0x0000\_0000} & \code{0xffff\_f000} & The boot memory \\
6 & \code{0x1000\_0000} & \code{0xffff\_0000} & The debug module \\
\bottomrule
\end{tabular}
\end{center}

\dstables{Router}

\subsection*{Behaviour}

\begin{itemize}
\item \textbf{Write address.} The head of \code{aw} is taken when no
write burst is running (\code{wbusy} clear) and the first peripheral
whose range matches has room on its \code{aws} port, or, for a hole,
when no write \code{DecErr} is owed. It is sent to that peripheral
alone.
\item Taking it sets \code{wbusy}, and records the peripheral in
\code{wsel}, a bit each, and the identifier in \code{wid}.
\code{whole} records that no range matched.
\item \textbf{Write beats.} While \code{wbusy} is set, beats pass from
\code{w} to the peripheral in \code{wsel} when it has room. A burst to
a hole has no bit in \code{wsel}, and its beats are taken and dropped.
The beat marked \code{last} clears \code{wbusy}.
\item So a second write address phase, to any peripheral, waits for
the first burst's last beat. AXI4 puts no identifier on \code{w}, and
this keeps two bursts' beats apart.
\item \textbf{Write to a hole.} When its last beat has gone, the router
sets \code{berr} and \code{berr\_id}. It sends a \code{B} of
\code{DecErr} with that identifier when \code{b} has room. That answer
goes before any peripheral's.
\item \textbf{Read address.} The head of \code{ar} is taken when the
first peripheral whose range matches has room on its \code{ars} port,
or, for a hole, when \code{rerr} is clear. Several reads may be
outstanding at once.
\item \textbf{Read to a hole.} Taking it sets \code{rerr}, keeps its
identifier in \code{rerr\_id}, and sets \code{rerr\_left} to its
\code{len} plus one. The router then sends that many beats of
\code{DecErr} and zero data on \code{r}, the last marked \code{last}.
\item \textbf{Read beats.} When no read burst is going up
(\code{rbusy} clear), the router's own error burst goes first. Else the
lowest peripheral that offers a beat is chosen.
\item A peripheral's beat that goes up without \code{last} sets
\code{rbusy} and records the peripheral in \code{rsel}. Its beats then
go up alone until its \code{last}, so \code{last} still marks the end
of one burst.
\item \textbf{Write responses.} When the router owes no \code{DecErr},
the lowest peripheral that offers a \code{B} sends it to the host. Each
response keeps the identifier the peripheral gave it.
\end{itemize}

\subsection*{Verification}

\begin{itemize}
\item Two tests in \code{bus::router} run \code{Router<2, TwoMap, 16,
32, 4, 3>} between a host tracker and two peripheral trackers, in
\code{//lib/parts:parts\_test}. \code{each\_burst\_reaches\_its\_range}
reads from both ranges and reads three beats from a hole, and checks
the hole's answer is \code{DecErr} in three beats.
\code{write\_beats\_do\_not\_interleave} writes three beats to each
peripheral, one burst after the other, and checks each peripheral got
its own burst whole and in order.
\item \code{ex\_axi\_router} issues five bursts before collecting any:
a read to each of two peripherals, a write to the third, and a read and
a write to a hole. The build simulates the netlist of \code{Router<3,
ThreeMap, 16, 32, 4, 3>} against that run under nvc and Verilator:
\code{//docs:axi\_router\_sim\_axi\_router\_tb\_test} and
\code{//docs:axi\_router\_vsim\_test}.
\item Vreteno's run lowers its router, \code{Router<3, RunMap, 32, 32,
4, 2>} in \code{cpu/vreteno/src/main.rs}, and the build simulates that
netlist against the run: \code{//docs:vreteno\_sim\_router\_router\_tb\_test}
and \code{//docs:vreteno\_vsim\_router\_test}.
\code{//cpu/vreteno:lockstep\_test} runs the same router in simulation.
\item \code{four\_peripherals\_share\_one\_node} in \code{bus::noc} puts
a \code{Router<4, NocMap, ..>} behind a node of the network, with a
memory in each range. Two hosts write and read every peripheral and
read a hole.
\end{itemize}

\subsection*{Used by}

\code{Router<3, RunMap, ..>} in Vreteno's run and its lockstep and
pair tests; \code{Router<4, RunMap, ..>}, with the boot memory, in
\code{cpu/vreteno/src/run.rs}; \code{Router<7, BoardMap, ..>} on the
board; \code{Router<4, NocMap, ..>} in the network's test;
\code{Router<3, ThreeMap, ..>} in \code{ex\_axi\_router}; and
\code{Router<2, TwoMap, ..>} in the router's own tests.

\subsection*{Limits}

It serves one host. Several hosts contending for one peripheral are not
handled here, as \code{//docs:axi} states. Write bursts pass one at a
time, whichever peripherals they go to. The merge of responses is by
fixed priority, lowest peripheral first. A mask compares only its own
bits: on an address wider than 16 bits, a mask of \code{0xf000} also
matches its base plus every multiple of \code{0x10000}. Nothing checks
that ranges do not overlap; an address in two ranges goes to the first.
The router decodes on the address alone, and ignores \code{region} and
\code{qos}.

% SPDX-License-Identifier: Apache-2.0
% covers: LiteBridge
\datasheet{LiteBridge}{The AXI4 to AXI-Lite bridge: one AXI4 link to any count of AXI-Lite peripherals, by an address map.}

\dsfacts{\code{lib/parts/src/bus/axi\_lite.rs}}%
{\code{txhdl\_parts::bus::axi\_lite::LiteBridge<N, M, ..>}, with \code{M} an address map, \code{txhdl::map::AddrMap<N>}}%
{\code{//docs:axi}, \code{//docs:vreteno}, \code{//docs:eth}, \code{//docs:hdmi}}

\subsection*{Function}

AXI-Lite is AXI4 with no identifier, no burst and no \code{last}, so
one transaction is one beat on each channel it uses. A bridge stands
between an AXI4 link and that many AXI-Lite peripherals. On its AXI4
side it takes the \code{aw}, \code{ar} and \code{w} an \code{AxiPer}
would take, and gives the \code{b} and \code{r} it would give. So it
goes where an \code{AxiPer} would: behind a host's tracker, or behind
one of a router's peripheral ports.

Its peripheral side is arrays of \code{N} ports: it drives
\code{aws}, \code{ars} and \code{ws}, and reads \code{bs} and
\code{rs}, entry \textit{i} of each being peripheral \textit{i}'s, so
the netlist's ports are \code{aws\_}\textit{i} and the like. An
address phase there is a \code{LiteAddr}: \code{addr} and
\code{prot}. A write word is a \code{LiteW}: \code{data} and
\code{strb}. \code{axi\_lite} makes the five channels of such a link.
One unit serves every count; the lowering unrolls its loops over the
arrays when it runs (issue 500). Until issue 647 a macro,
\code{lite\_bridge!}, wrote \code{LiteBridge1} to \code{LiteBridge8}
instead, one unit per count, with each range a pair of const
parameters.

\subsection*{Parameters}

\begin{center}\footnotesize
\begin{tabular}{@{}lp{0.7\textwidth}@{}}
\toprule
Parameter & Meaning \\
\midrule
\code{A} & The address width, in bits, on both sides. \\
\code{D} & The data width, in bits, on both sides. \\
\code{S} & The write strobe width, one bit per byte lane, so \code{D / 8}. \\
\code{I} & The identifier width of the AXI4 side, in bits. \\
\code{N} & The count of peripherals, one or more. \\
\code{M} & The address map: a type implementing \code{AddrMap<N>},
whose \code{RANGES} holds each peripheral's base and mask, in the
order of the ports. \\
\bottomrule
\end{tabular}
\end{center}

\subsection*{Address map}

An address \code{a} is peripheral \textit{i}'s when
\code{(a \& mask) == base} for \code{(base, mask) = M::RANGES[}\textit{i}\code{]}.
The bridge decodes a burst's first address, and the lowest range that
matches wins. An address in no range is a hole. The tables below are
generated for \code{LiteBridge<1, SerialMap, 32, 32, 4, 2>}, the bridge
in front of Vreteno's serial port, whose map is this one:

\begin{center}\footnotesize
\begin{tabular}{@{}llll@{}}
\toprule
Peripheral & \code{BASE} & \code{MASK} & In Vreteno \\
\midrule
0 & \code{0x3000} & \code{0xf000} & The serial port \\
\bottomrule
\end{tabular}
\end{center}

The board's design uses a bridge of six, \code{SlotMap}, at
\code{0x3000} to \code{0x3500} with mask \code{0xffff\_ff00} each, and
bridges of one at \code{0x0c00\_0000} (mask \code{0xfc00\_0000}, in
front of its interrupt controller) and at \code{0x1000\_0000} (mask
\code{0xffff\_0000}, the debug module). \code{ex\_axi\_lite} uses a
bridge of two with ranges at \code{0x1000} and \code{0x2000}, mask
\code{0xf000} each.

\dstables{LiteBridge}

\subsection*{Behaviour}

\begin{itemize}
\item \textbf{Taking a burst.} It takes a burst only when none is in
progress (\code{busy} clear). When both address channels offer,
\code{rfirst} says which goes first. After a read the write channel
goes first, and after a write the read channel.
\item Taking a burst records its identifier in \code{xid}, its address
in \code{addr}, its \code{len} in \code{left}, its \code{prot}, and
its peripheral in \code{sel}, a bit each. A hole has no bit set.
\code{sel} is \code{N} bits wide; a bridge of one has a one-bit
\code{sel}, which the netlist writes as a single bit.
\item \code{stride} is how far the address moves per beat:
\code{1 << size} bytes, or zero for a \code{Fixed} burst. Every other
burst kind moves as \code{Incr} does.
\item \textbf{A write beat.} The bridge takes the burst's next beat
from \code{w} when the chosen peripheral has room on both its
\code{aw} and its \code{w}. It sends that peripheral \code{addr} and
\code{prot} on \code{aw}, and the beat's \code{data} and \code{strb}
on \code{w}. To a hole, the beat is taken and dropped.
\item \textbf{A read beat.} The bridge sends \code{addr} and
\code{prot} on the chosen peripheral's \code{ar} when it has room.
\item Once a beat's request is out (\code{sent}), the bridge waits for
that peripheral's answer. A hole's answer is there at once:
\code{DecErr}, with zero data for a read.
\item \textbf{A read's answer} goes up on \code{r} as a beat of the
burst, when \code{r} has room. It has the burst's identifier, the word,
the peripheral's response, and \code{last} on the final beat.
\item \textbf{A write's answer} is taken when \code{b} has room, or when
more beats are to come. \code{wresp} keeps the first response that is
not \code{Okay}. After the last beat one \code{B} goes up, with the
burst's identifier and that response.
\item After each answer, \code{addr} moves on by \code{stride} and
\code{left} counts down. The answer to the last beat clears
\code{busy}.
\end{itemize}

\subsection*{Verification}

\begin{itemize}
\item \code{bursts\_cross\_the\_bridge\_a\_beat\_at\_a\_time} in
\code{bus::axi\_lite}, in \code{//lib/parts:parts\_test}, puts
\code{LiteBridge<2, TwoMap, 16, 32, 4, 2>} behind
a host tracker, with two simulated memories. It writes three words and
reads four, reads one word three times in a fixed burst, uses the
second peripheral, and sends a write and a read of two beats to a hole.
It checks every word and response, and counts ten transactions on the
first memory and two on the second.
\item \code{ex\_axi\_lite} runs the same kind of bursts against two
banks of registers written as hardware. The build simulates the
bridge's netlist against that run under nvc and Verilator:
\begin{itemize}
\item \code{//docs:axi\_lite\_sim\_axi\_lite\_bridge\_axi\_lite\_bridge\_tb\_test};
\item \code{//docs:axi\_lite\_vsim\_axi\_lite\_bridge\_test}.
\end{itemize}
\item Vreteno's run lowers its serial port's bridge, the
bridge of one of the tables, and the build simulates that netlist
against the run: \code{//docs:vreteno\_sim\_ubridge\_ubridge\_tb\_test}
and \code{//docs:vreteno\_vsim\_ubridge\_test}.
\item \code{//cpu/vreteno:lockstep\_test} and \code{//soc:demo\_test} run
a bridge of one in simulation.
\end{itemize}

\subsection*{Used by}

A bridge of two in the example \code{ex\_axi\_lite}. A bridge of one
in front of the serial port in Vreteno's run
(\code{cpu/vreteno/src/run.rs} and \code{cpu/vreteno/src/main.rs}) and
in \code{soc/src/lib.rs}; in the board's design, a bridge of six as its
field \code{puart} and bridges of one as \code{pplic} and \code{pdm}.
A bridge of one also puts each peripheral of the examples behind a
host's tracker: \code{ex\_eth}, \code{ex\_hdmi}, \code{ex\_sd} and the
rest.

\subsection*{Limits}

It serves one burst at a time, for all its peripherals together, since
AXI-Lite has no identifier to tell two apart. A burst's peripheral is
chosen on its first address, and every beat of the burst goes there.
It treats a \code{Wrap} burst as \code{Incr}. It counts a write's beats
from \code{len} and does not read \code{last} on \code{w}. Of the
address phase only \code{addr} and \code{prot} reach a peripheral;
\code{size} sets the stride, and \code{lock}, \code{cache},
\code{qos} and \code{region} are dropped.

% SPDX-License-Identifier: Apache-2.0
% covers: AxiWb
\datasheet{AxiWb}{A bridge from the peripheral end of an AXI4 link to the host lines of a pipelined Wishbone bus.}

\dsfacts{\code{lib/parts/src/bus/wb.rs}}{\code{txhdl\_parts::bus::wb::AxiWb}}%
{\code{//docs:axi}, \code{//docs:vreteno}}

\subsection*{Function}

A controller written elsewhere often speaks Wishbone rather than AXI.
\code{AxiWb} puts such a controller on a link. It is a peripheral
client written as hardware. On one side it holds the channel ends an
\code{AxiPer} gives a client: requests on \code{req} and write beats on
\code{wd} in, answers on \code{ans} and read beats on \code{rb} out.
On the other side it drives the host lines of a pipelined Wishbone
bus, \code{cyc}, \code{stb}, \code{we}, \code{adr}, \code{dat} and
\code{sel}, and reads \code{stall}, \code{ack} and the peripheral's
word, \code{rdat}.

It takes one single-beat burst at a time and puts it on the lines as
one request. The DDR3 controller of Vreteno's board is behind one.

\subsection*{Parameters}

\begin{center}\footnotesize
\begin{tabular}{@{}lp{0.7\textwidth}@{}}
\toprule
Parameter & Meaning \\
\midrule
\code{A} & The link's address width, in bits. \\
\code{I} & The link's identifier width, in bits. \\
\code{AW} & The width of the Wishbone word address, in bits. \\
\bottomrule
\end{tabular}
\end{center}

Words are 32 bits with four lanes, fixed in the source. The tables
below use \code{AxiWb<32, 2, 28>}. There \code{AW} is \code{ddr3::AW},
28 bits, the word address of the board's DDR3 memory.

\dstables{AxiWb}

\subsection*{Behaviour}

\code{stage} says where the request is. The source names five values.

\begin{itemize}
\item \textbf{0, waiting.} A request on \code{req} is taken when one
offers. The bridge keeps its identifier and its word address, and sets
\code{wsel} to all four lanes. A read goes to stage 2 at once, and a
write to stage 1.
\item \textbf{1, a write waiting for its beat.} The beat on \code{wd}
is taken when it offers. Its \code{data} goes to \code{wdat} and its
\code{strb} to \code{wsel}. Then stage 2.
\item \textbf{2, offered.} \code{cyc} and \code{stb} are high. A cycle
with \code{stall} low is the one the peripheral takes the request in,
and the stage moves to 3.
\item \textbf{3, taken.} \code{cyc} alone is high until \code{ack}.
An \code{ack} in the cycle of the take counts too. On \code{ack} the
stage moves to 4, and a read keeps the word on \code{rdat} in
\code{wdat}.
\item \textbf{4, answering.} A read sends one beat on \code{rb}: the
identifier, \code{wdat}, \code{Okay}, and \code{last} high. A write
sends an \code{Answer} of \code{Okay} on \code{ans}. Each waits for
room on its channel, then the stage returns to 0.
\item \code{we} is high unless the request is a read. \code{adr},
\code{dat} and \code{sel} are \code{wadr}, \code{wdat} and \code{wsel}.
Every line depends only on the bridge's registers, so nothing
combinational passes from the link to the lines or back.
\item The Wishbone address is a word address: the burst's byte address
shifted down by two and cut to \code{AW} bits. A peripheral decoded at
a region of the link's map sees offsets from the region's base, as
long as the base lies above those bits. A write at
\code{0x4000\_0010} lands in word 4.
\end{itemize}

\subsection*{Verification}

\begin{itemize}
\item Two tests in \code{bus::wb}, in \code{//lib/parts:parts\_test},
run the bridge on a link, behind \code{AxiHost} and \code{AxiPer},
against \code{WbMem}, a simulated Wishbone memory. The first,
\code{a\_read\_gets\_what\_a\_write\_put}, writes and reads one word.
The second, \code{slow\_answers\_and\_a\_calibration}, writes and
reads back 24 words at pseudorandom addresses, against a memory that
stalls for its first 40 cycles and answers after 5.
\item \code{ex\_axi\_wb} writes three words and reads them back, against
a memory that stalls for 10 cycles and answers after 2. The build
simulates the netlist of \code{AxiWb<32, 2, 28>} against that run under
nvc and Verilator: \code{//docs:axi\_wb\_sim\_axi\_wb\_tb\_test} and
\code{//docs:axi\_wb\_vsim\_test}.
\item In \code{//ddr3:ddr3\_test}, \code{words\_go\_in\_and\_come\_back}
runs \code{Ddr3Per}, this bridge with the controller's simulation
model behind it, and
\code{the\_controller\_is\_instantiated\_and\_not\_written} checks that
its netlist writes the bridge's module.
\item \code{//cpu/vreteno/board/sim:board\_test}, a manual target,
simulates the whole board under Vivado's simulator.
\end{itemize}

\subsection*{Used by}

\code{Ddr3Per} in \code{ddr3/src/lib.rs}, as its field \code{bridge},
and so the board's DDR3 memory at \code{0x4000\_0000}. The example
\code{ex\_axi\_wb}.

\subsection*{Limits}

It handles single-beat bursts only. It takes one write beat per
request, and answers a read with one beat marked \code{last}. It holds
one request at a time. It always answers \code{Okay}, and it reads no
error line from the peripheral. Words are 32 bits with four lanes.
Address bits above the \code{AW} bits it keeps are dropped.

% The network on chip.
% SPDX-License-Identifier: Apache-2.0
% covers: Switch
\datasheet{Switch}{A five-port switch of the network on chip, one per virtual channel.}

\dsfacts{\code{lib/parts/src/bus/noc/switch.rs}}{\code{txhdl\_parts::bus::noc::switch::Switch}}%
{\code{//docs:noc}}

\subsection*{Function}

A switch has a two-way link on each of its four sides and a fifth port, the exit,
where something that is not a node attaches. Its inputs are
\code{n\_in}, \code{s\_in}, \code{w\_in}, \code{e\_in} and
\code{x\_in}, and its outputs are \code{n\_out}, \code{s\_out},
\code{w\_out}, \code{e\_out} and \code{x\_out}. Each port passes one
\code{Pkt}, a beat of one AXI channel addressed to a node. The switch
routes in dimension order, X before Y. Its place in the lattice is two more
inputs, \code{col} and \code{row}, which whoever places it holds at
constants with \code{tie}, so every switch of a lattice is one type
(issue 635). A \code{Node} is two switches,
one per virtual channel.

\subsection*{Parameters}

\begin{center}\footnotesize
\begin{tabular}{@{}lp{0.7\textwidth}@{}}
\toprule
Parameter & Meaning \\
\midrule
\code{XB} & The width of a column coordinate, and of \code{col}, in
bits. A lattice is at
most \code{1 << XB} columns wide. \\
\code{YB} & The width of a row coordinate, and of \code{row}, in bits. A lattice is at
most \code{1 << YB} rows high. \\
\code{A} & The address width of the AXI link. \\
\code{D} & The data width of the AXI link. \\
\code{S} & The strobe width, which is \code{D / 8}. \\
\code{I} & The identifier width of the AXI link. \\
\bottomrule
\end{tabular}
\end{center}

The tables below are for \code{Switch<2, 2, 32, 32, 4, 2>}: two-bit
coordinates, on the link that \code{//soc} uses.

\dstables{Switch}

\subsection*{Behaviour}

\begin{itemize}
\item The switch holds one thing, \code{turn}, three bits: which
input the choice starts from. The port a packet leaves by holds
nothing, being a function of the packet's \code{dx} and \code{dy} and
of \code{col} and \code{row}; what needed a register is which input goes
when several could (issue 315).
\item The switch compares coordinates at \code{CW}, 8 bits. It sends a
packet east when \code{dx} is above \code{col}, and west when \code{dx}
is below it. In the right column, it sends the packet south when
\code{dy} is above \code{row}, and north when \code{dy} is below it. At its own node it sends the packet out of the exit.
\item Each cycle the switch works out, for every input with a packet,
the port that packet wants and whether that port is ready. An input
whose port is full does not hold up the others.
\item Of the inputs that can move, the switch takes one per cycle: the
first at or after \code{turn}, going round the five, and \code{turn}
becomes the one after whichever moved. So an input waits on at most
four others, where the fixed priority it replaced could starve one
outright: a busy link kept the exit from moving at all, since the exit
was last and the links were never all quiet.
\item The switch sends the packet it takes unchanged.
\end{itemize}

\subsection*{Verification}

\begin{itemize}
\item Two tests in \code{lib/parts/src/bus/noc.rs}, in
\code{//lib/parts:parts\_test}, drive one switch at 0,0 alone.
\code{a\_busy\_link}\allowbreak\code{\_and\_the\_exit\_take\_turns} has a link and the exit
offer every cycle for the east port, and asserts that each moves more
than fifty of two hundred and that the two differ by at most two.
\code{two\_inputs\_of\_a\_switch}\allowbreak\code{\_contend\_for\_their\_shared\_output}
measures two inputs leaving by one port, both offering every cycle and
one offering every third.
\item Eight more tests of the network in the same file run it inside
four nodes, and \code{a\_core\_draws\_a\_solid\_and\_a\_gpu\_fills\_it},
in \code{//soc:demo\_test}, runs a core, a GPU, a memory and a serial
port on a two by two lattice of nodes.
\item The switch's netlist is checked inside \code{Mesh}, two to a
node: \code{ex\_mesh} lowers a mesh of six nodes, and
\code{//docs:mesh\_sim\_mesh\_tb\_test} and \code{//docs:mesh\_vsim\_test}
run it under nvc and Verilator against the Rust run, at the mesh's
ports.
\end{itemize}

\subsection*{Used by}

\code{Node}, which holds two: \code{req} for requests and \code{rsp}
for responses.

\subsection*{Limits}

The switch moves one packet per cycle, not one per port.
\code{//docs:noc} measures the cost on the system of \code{//soc}: the
rasteriser draws its list in 36\,516 cycles on a direct link and in
78\,391 cycles across the lattice.

% SPDX-License-Identifier: Apache-2.0
% covers: Node
\datasheet{Node}{A node of the network on chip: a switch per virtual channel.}

\dsfacts{\code{lib/parts/src/bus/noc/node.rs}}{\code{txhdl\_parts::bus::noc::node::Node}}%
{\code{//docs:noc}}

\subsection*{Function}

A node is two \code{Switch} units: \code{req}, for what a host sends
and a peripheral receives, and \code{rsp}, for what a peripheral sends
and a host receives. The two share nothing. So a request held up behind
a full buffer cannot hold up the response that would empty it.

The node's place in the lattice is two inputs, \code{col} and
\code{row}, which it passes to both switches. Whoever places the node
holds them at constants with \code{tie}, so every node of a lattice is
one type (issue 635), and \code{//docs:noc} states that synthesis
folds the constants into the comparisons that route a packet.

Its other twenty ports are written out one by one. A port name starts
with \code{q} for the request channel or \code{p} for the response
channel. The side follows: \code{n}, \code{s}, \code{w}, \code{e}, or
\code{x} for the exit. Last is \code{i} or \code{o}. So \code{qei} is
the request input from the east, and \code{pxo} is the response output
to the exit.

\subsection*{Parameters}

\begin{center}\footnotesize
\begin{tabular}{@{}lp{0.7\textwidth}@{}}
\toprule
Parameter & Meaning \\
\midrule
\code{XB} & The width of a column coordinate, and of \code{col}, in
bits. A lattice is at most \code{1 << XB} columns wide. \\
\code{YB} & The width of a row coordinate, and of \code{row}, in bits.
A lattice is at most \code{1 << YB} rows high. \\
\code{A} & The address width of the AXI link. \\
\code{D} & The data width of the AXI link. \\
\code{S} & The strobe width, which is \code{D / 8}. \\
\code{I} & The identifier width of the AXI link. \\
\bottomrule
\end{tabular}
\end{center}

The tables below are for \code{Node<2, 2, 32, 32, 4, 2>}, the node on
the link that \code{//soc} uses.

\dstables{Node}

\subsection*{Behaviour}

\begin{itemize}
\item The node has no logic of its own. It runs \code{req} on the ten
\code{q} ports and \code{rsp} on the ten \code{p} ports, side by side,
both at the place \code{col} and \code{row} say. What it holds is its
two switches' state, which is six bits: a \code{turn} of three in
each.
\item Each switch routes in dimension order and moves one packet per
cycle, choosing among the inputs that can move by a round robin of its
own, as the \code{Switch} sheet states. The two switches take turns
independently: a busy request channel does not decide which response
moves.
\item A host's bridge sends requests into \code{qxi} and takes
responses from \code{pxo}. A peripheral's bridge takes requests from
\code{qxo} and sends responses into \code{pxi}.
\end{itemize}

\subsection*{Verification}

\begin{itemize}
\item Eight tests in \code{lib/parts/src/bus/noc.rs}, in
\code{//lib/parts:parts\_test}, run four nodes on a two by two lattice
with both bridges, both AXI trackers and memory. Among them,
\code{a\_burst\_crosses\_the\_lattice\_and\_its\_answer\_comes\_back}
writes a word from the host at 0,0 to a memory at 1,1 and reads it
back; \code{two\_hosts\_share\_a\_memory\_and\_each\_answer\_goes\_home}
puts hosts at 0,0 and 1,0 on that memory; and
\code{four\_peripherals\_share\_one\_node} puts a router of four
memories behind the node at 1,1. The others cross the lattice with
reads of four beats, writes of two and of sixteen, a fixed write and a
wrapping one.
\item \code{a\_core\_draws\_a\_solid\_and\_a\_gpu\_fills\_it}, in
\code{//soc:demo\_test}, runs the core, the GPU, a memory and a serial
port on the four corners of one lattice.
\item The node's netlist is checked inside \code{Mesh}: \code{ex\_mesh}
lowers a mesh of six nodes, and \code{//docs:mesh\_sim\_mesh\_tb\_test}
and \code{//docs:mesh\_vsim\_test} run it under nvc and Verilator
against the Rust run, at the mesh's ports.
\end{itemize}

\subsection*{Used by}

The tests of the network in \code{lib/parts/src/bus/noc.rs}, and
\code{soc::run} in \code{soc/src/lib.rs}, each with four nodes wired by
\code{mesh::lattice}; and \code{Mesh}, which holds \code{N} of them.

\subsection*{Limits}

\code{mesh::lattice} makes the channels in simulation and does not
lower; \code{Mesh} is the same wiring as a unit that does. On a lattice
that \code{mesh::lattice} makes, a link at the edge drives a channel
nobody reads. Each of the two switches moves one packet per cycle, not
one per port.

% SPDX-License-Identifier: Apache-2.0
% covers: HostBridge
\datasheet{HostBridge}{The host side of an exit: AXI from a host in, packets into the network out.}

\dsfacts{\code{lib/parts/src/bus/noc/bridge.rs}}{\code{txhdl\_parts::bus::noc::bridge::HostBridge}}%
{\code{//docs:noc}}

\subsection*{Function}

\code{HostBridge} stands where an AXI host's five channels would go to
a peripheral, and sends them across the network instead. It reads the
host tracker's \code{aw}, \code{ar} and \code{w}, and drives its
\code{b} and \code{r}. On the network side it sends request packets on
\code{req} into its node's exit, and takes response packets from
\code{rsp}. It addresses each request by a map of three ranges.

\subsection*{Parameters}

\begin{center}\footnotesize
\begin{tabular}{@{}lp{0.7\textwidth}@{}}
\toprule
Parameter & Meaning \\
\midrule
\code{X}, \code{Y} & The column and row of the bridge's node. The
bridge stamps them on every request as its source. \\
\code{XB}, \code{YB} & The widths of a column and a row coordinate, in
bits. \\
\code{A} & The address width of the AXI link. \\
\code{D} & The data width of the AXI link. \\
\code{S} & The strobe width, which is \code{D / 8}. \\
\code{I} & The identifier width of the AXI link. \\
\code{B0}, \code{M0}, \code{X0}, \code{Y0} & The first range. An
address whose bits under the mask \code{M0} equal \code{B0} goes to
the node at \code{X0}, \code{Y0}. \\
\code{B1}, \code{M1}, \code{X1}, \code{Y1} & The second range, in the
same way. The first range wins where both match. \\
\code{B2}, \code{M2}, \code{X2}, \code{Y2} & The third range, the
default route. The doc comment says to give it a mask of zero. The
bridge's \code{run} does not read \code{B2} or \code{M2}: it sends
every address the first two ranges miss to \code{X2}, \code{Y2}. \\
\bottomrule
\end{tabular}
\end{center}

\subsection*{Address map}

The tables below are for \code{soc::Bridge<0, 0>}, the bridge of the
core at 0,0 in \code{soc/src/lib.rs}. That type is
\code{HostBridge<0, 0, 2, 2, 32, 32, 4, 2, 0x3000, 0xffff\_f000, 1, 1,
0x1000, 0xffff\_f000, 0, 1, 0, 0, 0, 1>}. Both hosts of \code{//soc}
use this map.

\begin{center}\footnotesize
\begin{tabular}{@{}llllp{0.4\textwidth}@{}}
\toprule
Range & Base & Mask & Node & What is there \\
\midrule
0 & \code{0x3000} & \code{0xffff\_f000} & 1,1 & The serial port. \\
1 & \code{0x1000} & \code{0xffff\_f000} & 0,1 & The program's data, in
the memory. \\
2 & \code{0} & \code{0} & 0,1 & Everything else, the display list and
the framebuffer included, in the memory. \\
\bottomrule
\end{tabular}
\end{center}

The masks reach the whole address. \code{//docs:noc} states why: with
a mask of \code{0xf000}, the framebuffer's rows at \code{0x13000} and
above went to the serial port.

\dstables{HostBridge}

\subsection*{Behaviour}

\begin{itemize}
\item The bridge holds nothing between cycles. A write is one packet,
so it has no burst to keep track of.
\item A write goes when \code{aw} offers an address phase, \code{w}
offers a beat, and \code{req} has room. The bridge takes both and sends
one packet with \code{chan} \code{W}. The packet holds the address
phase from \code{aw}, the data and strobe from \code{w}, and
\code{last} high.
\item A read goes when no write goes in that cycle, \code{ar} offers an
address phase, and \code{req} has room. The bridge takes it and sends
one packet with \code{chan} \code{Ar}.
\item The bridge sends at most one request packet per cycle, and a
write before a read.
\item Every request packet has the destination the map gives for its
address, the source \code{X}, \code{Y}, and the host's own identifier.
\item A response packet with \code{chan} \code{B} goes to \code{b} as
its identifier and response, when \code{b} is ready. Any other response
packet goes to \code{r} as its identifier, data, response and
\code{last}, when \code{r} is ready.
\end{itemize}

\subsection*{Verification}

\begin{itemize}
\item Three tests in \code{lib/parts/src/bus/noc.rs}, in
\code{//lib/parts:parts\_test}, put host bridges on a two by two
lattice with 16-bit addresses. Their map sends \code{0x1000} and
\code{0x2000}, each under the mask \code{0xf000}, and the default, all
to the node at 1,1.
\item \code{a\_burst\_crosses\_the\_lattice\_and\_its\_answer\_comes\_back}
has one host, at 0,0.
\item \code{two\_hosts\_share\_a\_memory\_and\_each\_answer\_goes\_home}
has hosts at 0,0 and 1,0 on one memory.
\item \code{four\_peripherals\_share\_one\_node} has hosts at 0,0 and
1,0, and four memories behind a router at 1,1.
\item \code{a\_core\_draws\_a\_solid\_and\_a\_gpu\_fills\_it}, in
\code{//soc:demo\_test}, runs \code{soc::Bridge} at 0,0 for the core
and at 1,0 for the rasteriser.
\item No waveform in \code{docs/BUILD.bazel} lowers the bridge, so its
netlist has no co-simulation under nvc or Verilator.
\code{//docs:showcase} lists the network's netlists as not yet
simulated.
\end{itemize}

\subsection*{Used by}

\code{soc::Bridge} in \code{soc/src/lib.rs}, for the core at 0,0 and
the rasteriser at 1,0. The tests of the network in
\code{lib/parts/src/bus/noc.rs}.

\subsection*{Limits}

A write of one beat is one packet. A burst of more than one beat does
not cross the network: it would need a virtual channel per source, or
reassembly at the peripheral, and the bridge has neither. The map has
three ranges. The last one takes every address the others miss. So no
address is unknown to the network, and neither bridge has to answer for
one.

% SPDX-License-Identifier: Apache-2.0
% covers: PerBridge
\datasheet{PerBridge}{The peripheral side of an exit: request packets in, AXI to a peripheral out.}

\dsfacts{\code{lib/parts/src/bus/noc/bridge.rs}}{\code{txhdl\_parts::bus::noc::bridge::PerBridge}}%
{\code{//docs:noc}}

\subsection*{Function}

\code{PerBridge} is the other end of \code{HostBridge}. It takes
request packets on \code{req} from its node's exit and unpacks them
into \code{aw}, \code{ar} and \code{w} for a peripheral's tracker. It
takes the tracker's answers on \code{b} and \code{r}, and packs them on
\code{rsp} back to the node each request came from.

It also renames. An identifier belongs to the host that made it, so
two hosts may use the same one. The bridge gives every request a local
identifier of its own, and remembers against it the source node and
the identifier that node used. The answer comes back under the local
identifier, and the bridge sends the packet home under the original.

\subsection*{Parameters}

\begin{center}\footnotesize
\begin{tabular}{@{}lp{0.7\textwidth}@{}}
\toprule
Parameter & Meaning \\
\midrule
\code{X}, \code{Y} & The column and row of the bridge's node. The
bridge stamps them on every answer as its source. \\
\code{XB}, \code{YB} & The widths of a column and a row coordinate, in
bits. \\
\code{A} & The address width of the AXI link. \\
\code{D} & The data width of the AXI link. \\
\code{S} & The strobe width, which is \code{D / 8}. \\
\code{I} & The identifier width of the AXI link. \\
\code{NIDS} & How many local identifiers the bridge gives out, and so
how many bursts the peripheral has in flight at most. \\
\bottomrule
\end{tabular}
\end{center}

The tables below are for \code{PerBridge<0, 1, 2, 2, 32, 32, 4, 2,
4>}, the bridge of the memory at 0,1 in \code{soc/src/lib.rs}.

\dstables{PerBridge}

\subsection*{Behaviour}

\begin{itemize}
\item The local identifier on offer is \code{turn}. It is free when its
bit in \code{busy} is low.
\item A request packet with \code{chan} \code{W} is a write. The bridge
takes it when \code{turn} is free and both \code{aw} and \code{w} are
ready. It sends the address phase on \code{aw} under the local
identifier, and the data and strobe on \code{w} with \code{last} high.
\item The bridge treats any other request packet as a read. It takes it
when \code{turn} is free and \code{ar} is ready, and sends the address
phase on \code{ar} under the local identifier.
\item On taking a request, the bridge stores the source column in
\code{sx}, the source row in \code{sy} and the original identifier in
\code{oid}, each at the local identifier. It sets that identifier's
\code{busy} bit and adds one to \code{turn}.
\item A request waits while the identifier whose turn it is has not
been answered. The bridge does not look for another free identifier.
\item The bridge sends at most one answer packet per cycle, when
\code{rsp} has room, and a write response before a read beat. The
packet goes to the \code{sx} and \code{sy} stored under the answer's
identifier, under the \code{oid} stored there, with the source
\code{X}, \code{Y}.
\item A write response goes with \code{chan} \code{B}, its response,
and \code{last} high. A read beat goes with \code{chan} \code{R}, its
data, its response and its \code{last}.
\item The bridge clears an identifier's \code{busy} bit when it sends a
write response, or the last beat of a read, under that identifier.
\end{itemize}

\subsection*{Verification}

\begin{itemize}
\item Three tests in \code{lib/parts/src/bus/noc.rs}, in
\code{//lib/parts:parts\_test}, put a
\code{PerBridge<1, 1, 2, 2, 16, 32, 4, 2, 4>} in front of the memory
at 1,1.
\code{a\_burst\_crosses\_the\_lattice\_and\_its\_answer\_comes\_back}
writes and reads one word from one host.
\code{two\_hosts\_share\_a\_memory\_and\_each\_answer\_goes\_home} checks
that each answer reaches the host that asked. \code{//docs:noc} states
that this test failed before the renaming was written.
\code{four\_peripherals\_share\_one\_node} puts a router of four
memories behind the bridge, and checks each word and the router's
\code{DecErr} for an address none of them has.
\item \code{a\_core\_draws\_a\_solid\_and\_a\_gpu\_fills\_it}, in
\code{//soc:demo\_test}, runs a bridge at 0,1 for the memory and one at
1,1 for the serial port.
\item No waveform in \code{docs/BUILD.bazel} lowers the bridge, so its
netlist has no co-simulation under nvc or Verilator.
\code{//docs:showcase} lists the network's netlists as not yet
simulated.
\end{itemize}

\subsection*{Used by}

\code{soc::run} in \code{soc/src/lib.rs}: at 0,1 in front of the
memory's \code{AxiPer}, and at 1,1 in front of the serial port, behind
\code{LiteBridge1}. The tests of the network in
\code{lib/parts/src/bus/noc.rs}.

\subsection*{Limits}

A write request is one packet with one beat, so a burst of more than
one beat does not cross the network. At most \code{NIDS} bursts are in
flight at the peripheral. \code{turn} is a register of \code{I} bits
and \code{busy} has \code{NIDS} bits. Every use in the tree sets
\code{NIDS} to 4 with \code{I} of 2. The bridge sends one answer per
cycle.

% Peripherals for the board.
% SPDX-License-Identifier: Apache-2.0
% covers: EthTx
\datasheet{EthTx}{The transmit half of the Ethernet MAC: whole frames from a channel onto GMII.}

\dsfacts{\code{lib/parts/src/eth.rs}}{\code{txhdl\_parts::eth::EthTx}}%
{\code{//docs:eth}}

\subsection*{Function}

\code{EthTx} takes a frame's bytes from the channel \code{tx}, each an
\code{EthByte}: the byte and whether it is the frame's last. A frame
there is the destination address, the source address, the type or
length and the payload. The unit stores the whole frame, then sends it
on GMII: \code{txd}, a byte per cycle, and \code{tx\_en}. It adds the
preamble, the start of frame delimiter, padding up to sixty bytes and
the four bytes of the CRC-32 frame check sequence. The unit runs on one
clock, so on a board it can run on the transmit clock.

\subsection*{Parameters}

\code{EthTx} has no type parameters. Its frame memory holds
\code{FRAME\_MAX} bytes, 2048, a constant of the module.

\dstables{EthTx}

\subsection*{Behaviour}

The unit is two processes. The store takes the frame in; the sender is
the wire's protocol written as the sequence it is: wait for a whole
frame, the preamble and the delimiter, the frame's bytes, zeros up to
sixty, the check sequence and the gap, a byte a cycle, each a turn of
a loop. The lowering numbers the sender's waits into a state register
the unit does not declare, \code{at\_wait}, keeps its loops' counts in
\code{for0} to \code{for3}, and holds the two outputs between states in
\code{txd\_held} and \code{tx\_en\_held}; all are in the tables above.

\begin{itemize}
\item While no frame is whole, the store takes a byte from \code{tx}
whenever one is offered, and stores it in \code{frame} at \code{fill}.
The byte marked last sets \code{whole}.
\item A whole frame starts the sender: seven bytes of \code{0x55}, then
\code{0xd5}, with the CRC register set to \code{CRC\_INIT}, every bit
set. The frame follows, the register running over it.
\item A frame of fewer than sixty bytes is followed by zeros up to
sixty, with the register still running.
\item The check sequence is the register complemented, least
significant byte first.
\item \code{tx\_en} is high from the preamble to the end of the check
sequence. Twelve bytes of gap follow with \code{tx\_en} low. Then
\code{sent} is up for one cycle, which adds one to \code{frames} and
lets the store forget the frame and take the next.
\item The store takes no byte from \code{tx} while a frame is whole,
from its last byte stored to the end of its gap, so it holds the client
off. The client may offer bytes as slowly as it likes.
\item \code{crc\_byte} steps the register by a whole byte: the byte
added to it, then eight shifts with the polynomial \code{0xedb88320}
going in whenever the bit that leaves is set.
\end{itemize}

\subsection*{Verification}

\begin{itemize}
\item The tests below are in \code{lib/parts/src/eth.rs}, and run in
\code{//lib/parts:parts\_test}.
\item \code{the\_byte\_wide\_crc\_is\_the\_bitwise\_one} checks
\code{crc\_byte} against the register stepped a bit at a time, on every
byte from four registers.
\item \code{the\_crc\_of\_a\_known\_string} checks \code{crc32} against
the catalogue value and the residue.
\item \code{a\_frame\_goes\_out\_as\_the\_model\_says\_and\_comes\_back}
loops the transmitter into \code{EthRx}, and holds the bytes on the
wire to \code{wire\_bytes}, for frames of 20 and 100 bytes.
\item \code{frames\_at\_the\_minimum\_length} does the same for frames
of 59, 60 and 61 bytes.
\item \code{ex\_eth} runs the unit behind \code{EthLite} with the wire
looped back, and asserts that it sent three frames. The build simulates
the netlist of \code{EthTx} against that run under nvc and Verilator:
\code{//docs:eth\_sim\_eth\_tx\_eth\_tx\_tb\_test} and
\code{//docs:eth\_vsim\_eth\_tx\_test}.
\end{itemize}

\subsection*{Used by}

The example \code{ex\_eth}. The echo design for the AX7A200B:
\code{//eth:netlist} writes the Verilog of \code{eth\_tx}, and
\code{eth/board/eth\_echo.v} instantiates it on the 125~MHz transmit
clock. \code{//eth:echo\_synth} and \code{//eth:echo\_pnr} put that
design through Vivado to a bitstream. \code{//docs:eth} states that it
has not yet run on the board.

\subsection*{Limits}

The unit holds one frame. A frame longer than \code{FRAME\_MAX} bytes
keeps overwriting its last byte. Only gigabit works: at 100 or
10~Mbit a PHY sends a nibble per cycle, and the unit does not take a
byte over two cycles.

% SPDX-License-Identifier: Apache-2.0
% covers: EthRx
\datasheet{EthRx}{The receive half of the Ethernet MAC: frames off GMII, checked, onto a channel.}

\dsfacts{\code{lib/parts/src/eth.rs}}{\code{txhdl\_parts::eth::EthRx}}%
{\code{//docs:eth}}

\subsection*{Function}

\code{EthRx} takes a frame off GMII: \code{rxd}, a byte per cycle,
\code{rx\_dv} and \code{rx\_er}. It stores the whole frame and checks
its CRC-32 frame check sequence. Only then does it offer the frame's
bytes on the channel \code{rx}, each an \code{EthByte}, without the
preamble and the check sequence. It drops a frame whose check fails.
The unit runs on one clock, so on a board it can run on the receive
clock the PHY recovers from the line.

\subsection*{Parameters}

\code{EthRx} has no type parameters. Its frame memory holds
\code{FRAME\_MAX} bytes, 2048, a constant of the module.

\dstables{EthRx}

\subsection*{Behaviour}

\begin{itemize}
\item \code{phase} is hunting 0, receiving 1 and offering 2.
\item While hunting, \code{rx\_dv} high with \code{0xd5} on \code{rxd}
starts a frame. The unit clears \code{len} and \code{bad}, and sets the
CRC register to \code{CRC\_INIT}. It passes over \code{0x55}.
\item While hunting, any other byte with \code{rx\_dv} high is the
middle of a frame whose start the unit missed. The unit sets
\code{skip} and ignores the line until \code{rx\_dv} falls.
\item While receiving, the unit stores a byte per cycle for as long as
\code{rx\_dv} stays high. The CRC register runs over every byte, the
check sequence included. \code{bad} records any cycle with
\code{rx\_er} high.
\item When \code{rx\_dv} falls, the frame is good if \code{bad} is low,
the register holds \code{CRC\_RESIDUE}, and \code{len} is more than
four. A good frame goes to offering. The unit drops any other frame and
adds one to \code{dropped}.
\item While offering, the unit sends one byte on \code{rx} each cycle
\code{rx} is ready. It offers \code{len} less four bytes, and marks the
last. After the last it adds one to \code{frames} and hunts again.
\item A frame that arrives while the unit is offering is dropped:
the unit sets \code{skip} and adds one to \code{dropped}, once for that
frame.
\end{itemize}

\subsection*{Verification}

\begin{itemize}
\item The tests below are in \code{lib/parts/src/eth.rs}, and run in
\code{//lib/parts:parts\_test}. Each loops \code{EthTx} into the
receiver a cycle late.
\item \code{a\_frame\_goes\_out\_as\_the\_model\_says\_and\_comes\_back}
checks that a frame of 20 bytes comes back padded to sixty, and a frame
of 100 bytes comes back as it went.
\item \code{frames\_at\_the\_minimum\_length} does the same at 59, 60
and 61 bytes.
\item \code{a\_corrupted\_frame\_is\_dropped} flips a byte of the second
of three frames. It checks that the other two come back and that one
frame is counted dropped.
\item \code{a\_frame\_on\_a\_busy\_receiver\_is\_dropped} holds the
client off. It checks that the second frame is dropped and counted, and
that the first is intact.
\item \code{ex\_eth} runs the unit with one byte of the second frame
flipped on the wire, and asserts that the receiver dropped one frame.
The build simulates the netlist of \code{EthRx} against that run under
nvc and Verilator: \code{//docs:eth\_sim\_eth\_rx\_eth\_rx\_tb\_test}
and \code{//docs:eth\_vsim\_eth\_rx\_test}.
\end{itemize}

\subsection*{Used by}

The example \code{ex\_eth}. The echo design for the AX7A200B:
\code{//eth:netlist} writes the Verilog of \code{eth\_rx}, and
\code{eth/board/eth\_echo.v} instantiates it on the PHY's receive
clock. \code{eth/board/eth\_cdc.v} takes its bytes across to the
transmit clock. \code{//eth:echo\_synth} and \code{//eth:echo\_pnr} put
that design through Vivado to a bitstream. \code{//docs:eth} states
that it has not yet run on the board.

\subsection*{Limits}

The unit holds one frame. A client that has not yet taken a frame
loses the next one. \code{len} stops at 2047, so the unit stores any
further bytes of a longer frame in the last slot. Only gigabit works:
at 100 or 10~Mbit a PHY sends a nibble per cycle, and the unit does not
take a byte over two cycles. The unit reads \code{rx\_er} only while
receiving.

% SPDX-License-Identifier: Apache-2.0
% covers: EthLite
\datasheet{EthLite}{The Ethernet peripheral on AXI-Lite: three words, a byte per transaction.}

\dsfacts{\code{lib/parts/src/eth.rs}}{\code{txhdl\_parts::eth::EthLite}}%
{\code{//docs:eth}}

\subsection*{Function}

\code{EthLite} stands between an AXI-Lite host and the two halves of
the MAC. It takes bytes from AXI-Lite writes and sends them to
\code{EthTx} on \code{tx}. It answers AXI-Lite reads with the bytes
\code{EthRx} offers on \code{rx}. Its AXI-Lite side is one port,
\code{bus}, a \code{LitePort} with 32-bit addresses and words:
\code{bus\_aw}, \code{bus\_ar} and \code{bus\_w} in, and \code{bus\_b}
and \code{bus\_r} out. \code{irq} is high while a received byte waits.

\subsection*{Parameters}

\code{EthLite} has no type parameters.

\subsection*{Register map}

The unit decodes the word from address bits 3 to 2, so the words are at
offsets \code{0x0}, \code{0x4} and \code{0x8} of its range.
\code{ex\_eth} puts them at \code{0x1000}, \code{0x1004} and
\code{0x1008}.

\begin{center}\footnotesize
\begin{tabular}{@{}lllp{0.55\textwidth}@{}}
\toprule
Word & Offset & Access & Bits \\
\midrule
0, status & \code{0x0} & read & 0: a received byte waits. 1: the
transmitter has room for a byte. \\
1, transmit & \code{0x4} & write & 7 to 0: the byte. 8: high on a
frame's last byte. \\
2, receive & \code{0x8} & read & 7 to 0: the oldest received byte. 8:
high on a frame's last byte. 9: high when a byte was there. \\
\bottomrule
\end{tabular}
\end{center}

\dstables{EthLite}

\subsection*{Behaviour}

\begin{itemize}
\item A write goes when \code{bus\_aw} and \code{bus\_w} both offer
and \code{bus\_b} is ready. A write to the transmit word also waits
until \code{tx} is ready, so the unit answers it only when the
transmitter has taken the
byte. A client that writes faster than the transmitter takes is held
off rather than losing bytes.
\item A write to the transmit word sends bits 7 to 0 and bit 8 of the
written data on \code{tx}. The unit answers every write
\code{Okay}; a write to any other word does nothing else.
\item A read goes when \code{bus\_ar} offers and \code{bus\_r} is ready.
The unit
answers every read \code{Okay}. A read of word 1 or word 3 answers
zero.
\item A read of the receive word takes the waiting byte, if one waits,
and answers with bit 9 high. With none waiting it answers with bit 9
low and takes nothing.
\item \code{irq} is high while \code{rx} offers a byte.
\item The unit holds no state. It does not read the write strobe.
\end{itemize}

\subsection*{Verification}

\begin{itemize}
\item \code{EthLite} has no \code{\#[test]} of its own.
\item \code{ex\_eth} puts it behind \code{LiteBridge1} and an
\code{AxiHost}, in front of \code{EthTx} and \code{EthRx} on a looped
wire. The host sends three frames, each as one fixed burst to the
transmit word, and reads them back in fixed bursts of sixteen from the
receive word. The run asserts that the first comes back padded to
sixty bytes, the second, corrupted on the wire, does not come back, and
the third comes back as it went.
\item The build simulates the netlist of \code{EthLite} against that
run under nvc and Verilator:
\code{//docs:eth\_sim\_eth\_lite\_eth\_lite\_tb\_test} and
\code{//docs:eth\_vsim\_eth\_lite\_test}.
\end{itemize}

\subsection*{Used by}

The example \code{ex\_eth}. \code{//docs:eth} states that the
peripheral is not yet on Vreteno's board, and the echo design for the
AX7A200B does not use it.

\subsection*{Limits}

The unit moves one byte per transaction. The AXI-Lite bridge makes a
fixed burst one transaction per beat at the same address, so a host
sends a whole frame as one burst. The unit checks no more of the
address than the word: the bridge in front sends it only its own range.

% SPDX-License-Identifier: Apache-2.0
% covers: Hdmi
\datasheet{Hdmi}{The video peripheral for an HDMI encoder chip: a raster, a framebuffer, and AXI-Lite to paint it.}

\dsfacts{\code{lib/parts/src/hdmi.rs}}{\code{txhdl\_parts::hdmi::Hdmi}}%
{\code{//docs:hdmi}}

\subsection*{Function}

\code{Hdmi} drives the parallel inputs of an encoder chip such as the
SiI9134 on the Alinx AX7A200B. It is one unit on the pixel clock. It
counts the raster, reads its framebuffer under the beam, and drives
\code{rgb}, 24 bits with red on top, \code{hsync}, \code{vsync} and
\code{de} from registers. A host paints the framebuffer through the
unit's AXI-Lite side: \code{aw}, \code{ar} and \code{w} in, \code{b}
and \code{r} out, with 32-bit addresses and words.

\subsection*{Parameters}

\begin{center}\footnotesize
\begin{tabular}{@{}lp{0.7\textwidth}@{}}
\toprule
Parameter & Meaning \\
\midrule
\code{HV} & Visible columns. \\
\code{HFP} & Columns of horizontal front porch. \\
\code{HSW} & Columns of horizontal sync. \\
\code{HBP} & Columns of horizontal back porch. \\
\code{VV} & Visible rows. \\
\code{VFP} & Rows of vertical front porch. \\
\code{VSW} & Rows of vertical sync. \\
\code{VBP} & Rows of vertical back porch. \\
\code{SHIFT} & A framebuffer pixel covers \code{1 << SHIFT} by
\code{1 << SHIFT} pixels of the screen. \\
\bottomrule
\end{tabular}
\end{center}

The tables below are for \code{Hdmi<640, 16, 96, 48, 480, 10, 2, 33,
2>}. Those are the constants of \code{hdmi::vga}, 640 by 480 at 60~Hz
as VESA states it, with \code{SHIFT} of 2. That is the mode of
\code{hdmi/src/netlist.rs}, and its picture is 160 by 120.

\subsection*{Register map}

The unit decodes the word from address bits 3 to 2, so the words are at
offsets \code{0x0}, \code{0x4} and \code{0x8} of its range.
\code{ex\_hdmi} puts them at \code{0x1000}, \code{0x1004} and
\code{0x1008}.

\begin{center}\footnotesize
\begin{tabular}{@{}lllp{0.5\textwidth}@{}}
\toprule
Word & Offset & Access & Bits \\
\midrule
0, status & \code{0x0} & read & 0: high in vertical blanking. 31 to 16:
the count of frames shown. \\
1, cursor & \code{0x4} & read, write & 7 to 0: the column. 14 to 8: the
row. \\
2, pixel & \code{0x8} & write & 11 to 0: the colour, four bits each of
red, green and blue from the top. \\
\bottomrule
\end{tabular}
\end{center}

\dstables{Hdmi}

\subsection*{Behaviour}

\begin{itemize}
\item \code{hc} counts the column and \code{vc} the row, each over the
visible part and the blanking. \code{frames} goes up by one at the end
of each frame.
\item At each edge the unit reads the framebuffer at the pixel under
the beam into \code{px}. \code{//docs:hdmi} calls this the synchronous
read a block RAM has. The framebuffer address is the row,
\code{vc >> SHIFT}, in the top seven bits and the column,
\code{hc >> SHIFT}, in the low eight.
\item At the same edge the unit registers the syncs in \code{hs\_q} and
\code{vs\_q}, low inside the sync pulse, and the enable in
\code{de\_q}, high over the visible part. So they meet the pixel they
belong to.
\item \code{rgb} is \code{px} with each 4-bit colour widened to 8 bits
by repeating its bits, so 15 is full scale.
\item A write to the cursor word sets \code{cx} and \code{cy}.
\item A write to the pixel word puts the colour into the framebuffer at
the cursor, and moves the cursor to the next column. From the last
column it moves to the first column of the next row, and from the last
row to the first. A fixed burst to this word paints a run of pixels.
\item A write goes when \code{aw} and \code{w} both offer and \code{b}
is ready. A read goes when \code{ar} offers and \code{r} is ready. The
unit answers every transaction \code{Okay}. A read of word 2 or word 3
answers zero, and a write to word 0 or word 3 does nothing else.
\end{itemize}

\subsection*{Verification}

\begin{itemize}
\item \code{the\_raster\_is\_where\_the\_mode\_puts\_it}, in
\code{//lib/parts:parts\_test}, runs \code{Hdmi<8, 2, 3, 3, 6, 1, 2, 2,
1>} for two frames. It checks that exactly one lag of the outputs
behind the counters fits \code{de}, \code{hsync} and \code{vsync}
together.
\item \code{ex\_hdmi} runs \code{Hdmi<16, 2, 3, 3, 12, 1, 2, 2, 1>}
behind \code{LiteBridge1} and an \code{AxiHost}. The host sets the
cursor to the corner and paints all 48 pixels of the 8 by 6
framebuffer in one fixed burst. The run rebuilds the last whole frame
from the pixels \code{de} marks, and asserts every one of its 192
screen pixels against the colour painted.
\item The build simulates that unit's netlist against the run under
nvc and Verilator: \code{//docs:hdmi\_sim\_hdmi\_video\_hdmi\_video\_tb\_test}
and \code{//docs:hdmi\_vsim\_hdmi\_video\_test}. The co-simulated mode
is the example's, not the one in the tables above.
\end{itemize}

\subsection*{Used by}

The example \code{ex\_hdmi}. The demonstration for the AX7A200B:
\code{//hdmi:netlist} writes the Verilog of \code{hdmi\_video} in the
640 by 480 mode, with a test picture in the framebuffer's initial
values. \code{hdmi/board/hdmi\_demo.v} instantiates it on a 25.2~MHz
pixel clock, with its AXI-Lite inputs held idle.
\code{//hdmi:demo\_synth} and \code{//hdmi:demo\_pnr} put the design
through Vivado to a bitstream. \code{//docs:hdmi} states that it has
not yet run on the board.

\subsection*{Limits}

A framebuffer pixel is 12 bits, and the framebuffer holds
\code{FB\_WORDS}, $2^{15}$ words, so the picture is at most 256
columns by 128 rows. \code{hc} and \code{vc} are 12 bits wide. The mode
is fixed when the type is named. The whole unit runs on the pixel
clock, bus included. A host on another clock needs a crossing of the
five AXI-Lite channels, and nothing in the tree provides one. The
framebuffer cannot be read over the bus. On the board, \code{//docs:hdmi}
marks two things as unconfirmed: which of the chip's data inputs are
red, green and blue, and the I/O standard of the chip's pins.

% SPDX-License-Identifier: Apache-2.0
% covers: I2cInit
\datasheet{I2cInit}{The I2C master that resets and configures an HDMI encoder chip.}

\dsfacts{\code{lib/parts/src/hdmi.rs}}{\code{txhdl\_parts::hdmi::I2cInit}}%
{\code{//docs:hdmi}}

\subsection*{Function}

\code{I2cInit} holds the chip's reset low on \code{nreset}, releases
it, waits, and then writes a table of register values as I2C
transactions. The table is the SiI9134's configuration,
\code{SII9134\_WRITES}, with two entries. The unit drives the two
open-drain lines through \code{scl\_low} and \code{sda\_low}, and reads
the data line back on \code{sda\_in}. \code{done} rises when the table
is written, and \code{failed} is high when the chip did not acknowledge
a byte.

The unit has no bus. It writes this table, one transaction per entry.
\code{//docs:hdmi} and the source both state that no document the part
was written from gives the chip's registers. The values are
unconfirmed until checked against the SiI9134's documentation or the
board vendor's example.

\begin{center}\footnotesize
\begin{tabular}{@{}llllp{0.4\textwidth}@{}}
\toprule
Entry & Address & Register & Value & Meaning, unconfirmed \\
\midrule
0 & \code{0x72} & \code{0x08} & \code{0x35} & System control: out of
power down, with the input bus as the board wires it. \\
1 & \code{0x7a} & \code{0x2f} & \code{0x00} & The output as DVI: no
HDMI packets, only the picture. \\
\bottomrule
\end{tabular}
\end{center}

The address is the chip's I2C address as eight bits, write bit
included.

\subsection*{Parameters}

\begin{center}\footnotesize
\begin{tabular}{@{}lp{0.7\textwidth}@{}}
\toprule
Parameter & Meaning \\
\midrule
\code{DIV} & Cycles in a quarter of an I2C bit. A bit is four
quarters. \\
\code{HOLD} & Cycles the reset is held low, and cycles the unit waits
after releasing it. \\
\bottomrule
\end{tabular}
\end{center}

The tables below are for \code{I2cInit<63, 2520000>}, the master of
\code{hdmi/src/netlist.rs}. Its doc comment states that a quarter of 63
cycles at 25.2~MHz makes a bit at 100~kHz, and that 2\,520\,000 cycles
are a tenth of a second.

\dstables{I2cInit}

\subsection*{Behaviour}

The unit is two processes. The quarter clock counts \code{tick} from 0
to \code{DIV} less one and round again, every cycle, and drives
\code{done} from \code{written} and \code{failed} from \code{nak}. The
other is the whole job written as the sequence it is: the reset held,
the reset released, then a transaction per entry, a quarter of a bit
per wait. The lowering numbers its waits into a state register the unit
does not declare, \code{at\_wait}, keeps its loops' counts in
\code{for0} to \code{for5}, and holds the three lines it drives
between states in \code{nreset\_held}, \code{scl\_low\_held} and
\code{sda\_low\_held}; all are in the tables above. \code{byte} is
\code{byte\_rw} in the netlist, since SystemVerilog reserves the word.

\begin{itemize}
\item \code{nreset} is low for \code{HOLD} cycles, then high, and the
unit waits \code{HOLD} cycles more before the first transaction.
\item A transaction is the start, the address, the register and the
value, each a byte followed by its acknowledge, and the stop, every
step a whole number of quarters of \code{DIV} cycles.
\item At the start the data line falls while the clock is high, then
the clock falls. At the stop the clock rises while the data line is
low, then the data line rises.
\item In a data bit the data line is set while the clock is low, and
the clock is high through the two middle quarters, so the data changes
only while the clock is low. Bytes go out most significant bit first,
from \code{shift}.
\item In an acknowledge the unit releases the data line, and reads
\code{sda\_in} at the end of the clock's high half, as it pulls the
clock low again. A high line sets \code{nak}, which \code{failed}
follows. Nothing clears \code{nak}.
\item \code{scl\_low} or \code{sda\_low} high pulls its line low.
\item After the stop of the last entry, \code{written} is set,
\code{done} rises a cycle later, and the sequence waits for ever.
\end{itemize}

\subsection*{Verification}

\begin{itemize}
\item \code{the\_chip\_is\_configured\_over\_i2c}, in
\code{//lib/parts:parts\_test}, runs \code{I2cInit<3, 5>} against
\code{I2cDevice}, a model of a device that acknowledges every byte and
records what it was sent. It checks that \code{done} rises, that
\code{failed} stays low, that the reset is released at cycle 5, and
that the device received \code{SII9134\_WRITES}.
\item \code{ex\_hdmi} runs \code{I2cInit<4, 10>} against the same model
and asserts the same table, \code{done}, and \code{failed} low.
\item The build simulates that unit's netlist against the run under
nvc and Verilator: \code{//docs:hdmi\_sim\_hdmi\_i2c\_hdmi\_i2c\_tb\_test}
and \code{//docs:hdmi\_vsim\_hdmi\_i2c\_test}.
\end{itemize}

\subsection*{Used by}

The example \code{ex\_hdmi}. The demonstration for the AX7A200B:
\code{//hdmi:netlist} writes the Verilog of \code{hdmi\_i2c}, and
\code{hdmi/board/hdmi\_demo.v} instantiates it with open-drain pads on
the chip's I2C lines and \code{done} and \code{failed} on LEDs.
\code{//hdmi:demo\_synth} and \code{//hdmi:demo\_pnr} put the design
through Vivado to a bitstream. \code{//docs:hdmi} states that it has
not yet run on the board.

\subsection*{Limits}

There is one table, and its values are unconfirmed. The unit keeps
writing after a byte is not acknowledged, and it writes the table once.
It has no input for the clock line. \code{tick} is 16 bits wide, so
\code{DIV} is at most 65\,536; the counters of the two holds are as
wide as \code{HOLD} needs, 22 bits in the tables.

% Vreteno and its devices.
% SPDX-License-Identifier: Apache-2.0
% covers: Vreteno
\datasheet{Vreteno}{A three-stage RV32IMC core with machine-mode traps and an AXI host port.}

\dsfacts{\code{cpu/vreteno/src/core.rs}}{\code{vreteno32::core::Vreteno}}%
{\code{//docs:vreteno}, \code{//docs:noc}, \code{//docs:tapeout}}

\subsection*{Function}

\code{Vreteno} is a RISC-V core for the RV32I base set, the M
extension and the C extension's compressed instructions. It is one
process with three stages. The fetch stage reads the instruction at
the program counter into the instruction register, expanding a
compressed one into the thirty-two bit instruction it stands for. The
execute stage decodes that word, computes, and issues
loads and stores on the bus. The writeback stage extends a loaded
word, writes the register file and retires the instruction.

The core holds two memories: the register file, 32 words, and the
instruction memory, \code{IMEM\_WORDS} or 1024 words.
\code{Vreteno::with} loads a program into the instruction memory. The
data memory and every device are on the bus, outside the core.

The bus side is an AXI host written as hardware. The core issues a
burst on \code{issue}, sends a store's beat on \code{wbeat}, and hands
identifiers back on \code{release}. It receives \code{grant},
\code{done} and \code{rdata} from its tracker, an \code{AxiHost}. The
other inputs are \code{rst}, the external interrupt \code{irq}, the
timer's line \code{tirq} and the software interrupt \code{sirq}, which
the controller raises from \code{msip}; and the debug module's, the
halt and resume requests \code{haltreq} and \code{resumereq} and its
register access, \code{dbg\_regno}, \code{dbg\_wdata} and
\code{dbg\_we}, honoured in debug mode (issue 154). The outputs are
\code{halt}, \code{instr}, the word retiring, \code{wb}, a
\code{Writeback} of \code{done}, \code{rd} and \code{val}, \code{dbg},
debug mode, and \code{dbg\_rdata}, the register the module asked for.

\subsection*{Parameters}

\begin{center}\footnotesize
\begin{tabular}{@{}lp{0.7\textwidth}@{}}
\toprule
Parameter & Meaning \\
\midrule
\code{IW} & The width of the AXI identifier on \code{rdata}, \code{done},
\code{grant} and \code{release}, in bits. The tables use 2, as the
board and the runs do. \\
\bottomrule
\end{tabular}
\end{center}

\subsection*{Register map}

The control registers, by CSR number. Software reaches them through
the six CSR instructions. They are not on the bus.

\begin{center}\footnotesize
\begin{tabular}{@{}llp{0.6\textwidth}@{}}
\toprule
Number & Register & What the core keeps \\
\midrule
\code{0x300} & \code{mstatus} & Bits 3 and 7 only; a write is masked
with \code{0x88}. \\
\code{0x304} & \code{mie} & Bits 11 and 7 only, the external and the
timer enable. \\
\code{0x305} & \code{mtvec} & The handler's address; a write clears
bits 1 and 0. \\
\code{0x340} & \code{mscratch} & The whole word. \\
\code{0x341} & \code{mepc} & The trap's address; a write clears
bit 0. \\
\code{0x342} & \code{mcause} & The whole word. \\
\code{0x343} & \code{mtval} & The whole word. \\
\code{0x344} & \code{mip} & Bit 11 is stored: \code{irq} sets it and
a write keeps only that bit. Bit 7 reads as \code{tirq}. \\
\code{0x7c0} & \code{mhalt} & Nothing: it reads as zero. A write of an
odd value stops the core when that instruction retires. \\
\code{0x301} & \code{misa} & Nothing: it reads
\code{0x4000\_1104}, which is \code{MXL} of 1 and the letters
\code{C}, \code{I} and \code{M}. Writable and ignored, as the
specification allows. \\
\code{0xf11} & \code{mvendorid} & Nothing: zero, which means
unimplemented. Read only. \\
\code{0xf12} & \code{marchid} & Nothing: zero. Read only. \\
\code{0xf13} & \code{mimpid} & Nothing: zero. Read only. \\
\code{0xf14} & \code{mhartid} & Nothing: zero, this machine's one
hart. Read only. \\
\bottomrule
\end{tabular}
\end{center}

A CSR instruction that names any other number is an illegal word and
traps.

\dstables{Vreteno}

\subsection*{Behaviour}

\begin{itemize}
\item With \code{rst} high, the fetch counter goes to zero, the
instruction register is emptied, and the load wait and the sequencer
are cleared.
\item The fetch reads the instruction memory at bits 11 to 2 of the
program counter and at the word after. The halfword at the counter is
the upper half of its word when bit 1 is set. If its low two bits are
not both set it is a compressed instruction, expanded in the fetch;
otherwise the halfword after it is the instruction's upper half. The
counter moves on by two or four.
\item \code{ir\_c} says the instruction in execute was compressed.
\code{jal} and \code{jalr} link the address two bytes on for a
compressed one and four for the rest.
\item An instruction in execute that reads the register the writeback
stage is about to write gets that value directly. When that value is a
load's, the instruction waits one cycle and reads the register file
instead.
\item A taken branch, \code{jal}, \code{jalr}, a trap and \code{mret}
redirect the fetch. The word fetched in that cycle is squashed, which
costs one bubble.
\item A load or store whose address has any bit set above bit 11 goes
on the bus as a burst of one beat, size 2, \code{INCR}. A store's beat
has the lanes it covers as its strobe.
\item A load or store waits for room on both \code{issue} and
\code{wbeat}, whatever its address.
\item A store is posted and the core runs on. A load sets
\code{dev\_wait} and sits in writeback, holding the pipeline, until
its answer lands in \code{wb\_dev}.
\item When a write response and a read beat arrive together, the
write's identifier goes back first and the read's waits a cycle. The
core drops what it receives on \code{grant}.
\item \code{ecall} traps with cause 11 and \code{ebreak} with cause 3,
which takes its own address as \code{mtval}. A word the core does not
know traps with cause 2, and \code{mtval} takes the word. A halfword
that is no RV32C instruction expands to itself, which no thirty-two bit
instruction is, so it traps with cause 2 and \code{mtval} takes the
halfword. An unaligned load traps with cause 4 and an unaligned store
with cause 6, and \code{mtval} takes the address that was asked for.
Every other trap writes zero to \code{mtval}.
\item On a trap, \code{mepc} takes the instruction's address, bit 7 of
\code{mstatus} takes bit 3, bit 3 is cleared, and the fetch goes to
\code{mtvec}. \code{mret} sets bit 3 from bit 7, sets bit 7, and goes
to \code{mepc}.
\item The external interrupt is ready when bit 11 is set in both
\code{mie} and \code{mip}. The software interrupt is ready when bit 3
of \code{mie} is set and \code{sirq} is high, and the timer's when bit
7 of \code{mie} is set and \code{tirq} is high. Any of the three is
taken when bit 3 of \code{mstatus} is set, in place of the instruction
in execute. The external interrupt, cause \code{0x8000\_000b}, wins
over the software interrupt, cause \code{0x8000\_0003}, which wins over
the timer's, cause \code{0x8000\_0007}.
\item A multiply or divide starts the sequencer and stalls in execute.
A multiply is one step and a division thirty-two. \code{//docs:vreteno}
states that a multiply takes three cycles and a division thirty-three.
An interrupt cancels the sequencer, and the instruction starts it
again after the handler returns.
\item The sequencer does not start while a device load in writeback is
still waiting for its answer.
\item A write of an odd value to \code{mhalt} stops the fetch and
parks the program counter one word past that instruction, which
retires first. \code{halt} rises as it retires, and stays high.
\item \code{stall}, \code{int\_take} and \code{redirect} are wires kept
as fields, so the trace and the netlist show them.
\end{itemize}

\subsection*{Verification}

\begin{itemize}
\item \code{//cpu/vreteno:lockstep\_test} runs the core with the router, the
data memory, the timer and the serial port against the reference model
in \code{model.rs}. After every cycle the program counter, registers
1 to 31, the eight control registers that hold something and the halt
must agree. At the
halt the data memory, the timer's compare and the serial port's bytes
must agree. Its tests are \code{demo\_program},
\code{a\_multiply\_right\_after\_a\_load},
\code{compressed\_instructions}, \code{an\_unaligned\_access\_traps},
and \code{random\_programs}, which
runs 64 seeds of 200 instructions each with both lengths mixed and
counts that the mix is there.
\item \code{//cpu/vreteno:isa\_test} compares the core's expander with
the decoder's on every compressed halfword, and
\code{//cpu/vreteno:compressed\_test} checks the decoder against the
GNU disassembler on all 49\,152 of them.
\item \code{//cpu/vreteno:hello\_test} (\code{rust\_runs\_on\_vreteno})
and \code{//cpu/vreteno:hello\_cc\_test}
(\code{cpp\_runs\_on\_vreteno}) run compiled programs on the machine
of \code{run.rs} and check what the serial line says and that the core
halted.
\item \code{//cpu/vreteno:board\_test} and \code{//soc:demo\_test} run
the core inside \code{Board} and on the lattice.
\item The build simulates \code{Vreteno<2>}'s netlist, with the
demonstration program in its instruction memory, against the trace of
\code{//cpu/vreteno:vreteno} under nvc and Verilator:
\code{//docs:vreteno\_sim\_vreteno\_vreteno\_tb\_test} and
\code{//docs:vreteno\_vsim\_vreteno\_test}.
\item \code{//cpu/vreteno:vreteno\_synth} and
\code{//cpu/vreteno:vreteno\_pnr} synthesise and place the core on an
\code{xc7a200tfbg484-2}; both are manual. \code{//docs:vreteno}
reports 2151 LUTs as logic, 48 LUTs as distributed RAM, 489
flip-flops, two block RAM tiles and four DSP blocks after synthesis. It
reports that the routed core meets a 10~ns clock with 0.296~ns of
setup slack.
\item \code{//cpu/vreteno/asic:vreteno\_cells} maps the core onto the
Nangate45 cells and builds with everything.
\code{//cpu/vreteno/asic:}\allowbreak\code{vreteno\_pnr} places and routes it, and is
manual.
\end{itemize}

\subsection*{Used by}

\code{Board}, as \code{cpu}. The demonstration
\code{//cpu/vreteno:vreteno}, the lockstep test, and \code{run.rs},
which the \code{hello} and \code{hello\_cc} runs call. The system in
\code{soc}, where the core is the host at corner 0,0 of the lattice.

\subsection*{Limits}

A program is at most 1024 words, 4096 bytes. The fetch uses bits 11
to 2 of the program counter, and the word after wraps too, so a
program that runs off the end wraps into itself.
Nothing loads the instruction memory at run time; a netlist gets its
contents through \code{init}.

A load from an address below \code{0x1000} sends no request. The
instruction retires with the word the previous device load left in
\code{wb\_dev}. A store below \code{0x1000} writes nothing.

A trap handler must start at a whole word: \code{mtvec} keeps no
bits 1 and 0. The core has machine mode only and the fourteen control
registers above, in three kinds: eight it keeps, \code{mhalt}, which
is its own, and five that say what the machine is and hold nothing. A
write to one of the four whose number begins \code{0xf} is an illegal
instruction, which is what read only means in the specification;
\code{misa} is writable and ignores what is written. The two counters answer as well:
\code{mcycle} at \code{0xb00} with its high half at \code{0xb80},
and \code{minstret} at \code{0xb02} with \code{0xb82}, each 64 bits
and each writable. \code{mcycle} counts while the core runs and stops
when the halt stops it; \code{minstret} counts what retires. They are
the one part of the core the lockstep test leaves alone, for the
reason \code{//docs:vreteno} gives, and
\code{//cpu/vreteno:counters\_test} checks them instead. It has one load in
flight at a time. On the board, \code{irq} is the interrupt
controller's line.

% SPDX-License-Identifier: Apache-2.0
% covers: Dmem
\datasheet{Dmem}{Vreteno's data memory: 1024 words behind an AXI peripheral end.}

\dsfacts{\code{cpu/vreteno/src/dmem.rs}}{\code{vreteno32::dmem::Dmem}}%
{\code{//docs:vreteno}}

\subsection*{Function}

\code{Dmem} is the data memory of the Vreteno machine, a device on its
bus. It holds \code{DMEM\_WORDS}, 1024 words, in four memories of a
byte, \code{lane0} to \code{lane3}, one per lane of the word. A store
of a byte or a half writes only its lanes and reads nothing.

It is the peripheral client of an \code{AxiPer}, and its link is one
port, \code{bus}, a \code{PerPort} of the four channels: it takes
requests on \code{bus\_req} and write beats on \code{bus\_w}, and
answers on \code{bus\_ans} for a write and \code{bus\_r} for a read.
It answers single-beat bursts,
which is all the core makes.

\subsection*{Parameters}

\begin{center}\footnotesize
\begin{tabular}{@{}lp{0.7\textwidth}@{}}
\toprule
Parameter & Meaning \\
\midrule
\code{I} & The width of the AXI identifier on \code{bus\_req},
\code{bus\_ans} and \code{bus\_r}, in bits. The tables use 4, as the board
does since its bus has two hosts behind an arbiter (issue 241). \\
\bottomrule
\end{tabular}
\end{center}

\subsection*{Address map}

The word a burst names is bits 11 to 2 of its address. Lane $k$ is
byte $k$ of the word: \code{lane0} is bits 7 to 0, \code{lane3} bits 31
to 24. The memory looks at no other address bits.

On the board, \code{BoardRouter} sends it the page at \code{0x1000}
with mask \code{0xffff\_f000}, so addresses \code{0x1000} to
\code{0x1ffc} are words 0 to 1023. The model's \code{DATA\_BASE} is
the same \code{0x1000}.

\dstables{Dmem}

\subsection*{Behaviour}

\begin{itemize}
\item A write is taken when no other write is waiting for its beat.
Its word and identifier go into \code{paddr} and \code{pid}, and
\code{pend} is set.
\item The held write's beat is taken when it has arrived and
\code{bus\_ans} has room. Each lane whose strobe bit is set takes its byte
of the beat, and \code{bus\_ans} sends \code{OKAY} with \code{pid}.
\item A read is taken when no write is waiting and the answer register
is free or is being emptied in this cycle. The four lanes are read
into \code{word} at the edge, the synchronous read a block RAM has.
\item The read is answered on \code{bus\_r} the cycle after, with
\code{rid}, \code{OKAY} and \code{last} set.
\item While a write waits for its beat, no request of either kind is
taken.
\item \code{Dmem::with} fills the lanes from bytes before the first
cycle. \code{data\_word} reads a word back for a run or a test.
\end{itemize}

\subsection*{Verification}

\begin{itemize}
\item \code{//cpu/vreteno:lockstep\_test} compares every word of the model's
data memory with the four lanes at the halt, in \code{demo\_program},
\code{a\_multiply\_right\_after\_a\_load} and \code{random\_programs}.
\item \code{//cpu/vreteno:hello\_test}, \code{//cpu/vreteno:hello\_cc\_test}
and \code{//cpu/vreteno:board\_test} put the compiled programs'
constants in it with \code{Dmem::with}, and check what the programs
print.
\item The build simulates \code{Dmem<2>}'s netlist against the trace
of \code{//cpu/vreteno:vreteno} under nvc and Verilator:
\code{//docs:vreteno\_sim\_dmem}\allowbreak\code{\_dmem\_tb\_test} and
\code{//docs:vreteno\_vsim\_dmem\_test}.
\end{itemize}

\subsection*{Used by}

\code{Board}, as \code{dmem} behind the tracker \code{pdmem}. The
demonstration \code{//cpu/vreteno:vreteno}, the lockstep test and
\code{run.rs}. \code{board\_netlist} writes a program's data into
\code{lane0} to \code{lane3} of the board's netlist.

\subsection*{Limits}

It serves bursts of one beat and does not look at a burst's length. It
checks no address, since the router sends it only its own range, so
an address beyond word 1023 wraps. One write is held at a time. Nothing
on the machine fills it at run time except stores, so a program's
constants have to be in it before the first cycle.
\code{cpu/vreteno/board/ax7a200.xdc} constrains the four lanes to block
RAM.

% SPDX-License-Identifier: Apache-2.0
% covers: Timer
\datasheet{Timer}{Vreteno's timer: a count, a compare and an interrupt line, behind an AXI peripheral end.}

\dsfacts{\code{cpu/vreteno/src/timer.rs}}{\code{vreteno32::timer::Timer}}%
{\code{//docs:vreteno}, \code{//docs:axi}}

\subsection*{Function}

\code{Timer} is the machine timer of the Vreteno machine. It holds a
64-bit count, \code{mtime}, and a 64-bit compare, \code{mtimecmp}, and
shows each as two words on the bus. It raises \code{tirq} while the
count has reached the compare. The core reads that line as bit 7 of
\code{mip}.

It is the peripheral client of an \code{AxiPer}: requests come on
\code{req} and write beats on \code{wd}, and answers go out on
\code{ans} and \code{rb}. The input \code{rst} clears the count.

\subsection*{Parameters}

\begin{center}\footnotesize
\begin{tabular}{@{}lp{0.7\textwidth}@{}}
\toprule
Parameter & Meaning \\
\midrule
\code{I} & The width of the AXI identifier on \code{req}, \code{ans}
and \code{rb}, in bits. The tables use 2, as the board does. \\
\bottomrule
\end{tabular}
\end{center}

\subsection*{Register map}

Offsets from the timer's base. On the board the base is
\code{0x2000}, which is also \code{isa::TIMER\_BASE}.

\begin{center}\footnotesize
\begin{tabular}{@{}llp{0.55\textwidth}@{}}
\toprule
Offset & Word & Access \\
\midrule
\code{0x0} & \code{mtime}, bits 31 to 0 & Read and write. \\
\code{0x4} & \code{mtime}, bits 63 to 32 & Read and write. \\
\code{0x8} & \code{mtimecmp}, bits 31 to 0 & Read and write. \\
\code{0xc} & \code{mtimecmp}, bits 63 to 32 & Read and write. \\
\bottomrule
\end{tabular}
\end{center}

The timer selects the word by bits 3 and 2 of the address and looks at
no other bits.

\dstables{Timer}

\subsection*{Behaviour}

\begin{itemize}
\item The count adds one every cycle. With \code{rst} high it goes to
zero.
\item A read is taken when \code{rb} has room and no write is waiting
for its beat. It is answered on \code{rb} in the cycle it is taken,
with the word addressed, \code{OKAY} and \code{last} set.
\item A write is taken when no write is waiting. The word it names
goes into \code{psel} and its identifier into \code{pid}.
\item The held write's beat is taken when it has arrived and
\code{ans} has room. The lanes the strobe covers replace those bytes of
the word, the other bytes stay, and \code{ans} sends \code{OKAY}.
\item Every cycle \code{pending} takes whether \code{mtime} is at or
above \code{mtimecmp}, compared as unsigned numbers. \code{tirq} is
\code{pending}, a register's output, so the core sees in one cycle
what the timer decided in the cycle before.
\item \code{mtimecmp} starts at zero and \code{rst} does not clear it,
so the line is high until software writes a compare above the count.
\end{itemize}

\subsection*{Verification}

\begin{itemize}
\item \code{//cpu/vreteno:lockstep} hands \code{pending} to the model
for each instruction and checks \code{mtimecmp} against the model at
the halt. \code{demo\_program} sets the compare to 150 and counts the
timer's interrupt. \code{random\_programs} writes the count and the
compare with arbitrary words and runs with the timer firing.
\item The build simulates \code{Timer<2>}'s netlist against the trace
of \code{//cpu/vreteno:vreteno} under nvc and Verilator:
\code{//docs:vreteno\_sim\_timer\_timer\_tb\_test} and
\code{//docs:vreteno\_vsim\_timer\_test}.
\end{itemize}

\subsection*{Used by}

\code{Board}, as \code{timer} behind the tracker \code{ptimer}, with
its line wired to the core's \code{tirq}. The demonstration
\code{//cpu/vreteno:vreteno}, the lockstep test and \code{run.rs}.

\subsection*{Limits}

The count counts clock cycles. The timer serves bursts of one beat and
checks no address beyond bits 3 and 2, since the router sends it only
its own range. One write is held at a time, and no read is taken while
it waits. A write updates 32 bits of a 64-bit register, so software
sets a whole compare in two stores.

% SPDX-License-Identifier: Apache-2.0
% covers: Plic1, Plic2, Plic3, Plic4, Plic5, Plic6, Plic7, Plic8
\datasheet{Plic}{Platform-level interrupt controllers: one to eight sources and one target, at the RISC-V PLIC's offsets.}

\dsfacts{\code{lib/parts/src/plic.rs}; generator \code{plic\_text} in \code{lib/macros/lib.rs}}%
{\code{txhdl\_parts::plic::Plic1} to \code{Plic8}}%
{\code{//docs:vreteno}}

\subsection*{Function}

A controller gathers the interrupt lines of several sources into one
line for one target, a hart's machine mode, and lets software ask which
source to serve. Each source \textit{i} is an input line,
\code{src}\textit{i}, numbered from 1. The output \code{irq} is the
target's line. Software reaches the controller as an AXI-Lite
peripheral: reads come on \code{ar} and are answered on \code{r}, and
writes come on \code{aw} and \code{w} and are answered on \code{b}. It
sits behind a \code{LiteBridge1}, which sends it only its own range.
The input \code{rst} clears the requests, the line, the priorities, the
enable bits and the threshold.

\code{plic!(PlicN, N)} writes one unit per count, since the lowering
reads a body and not a loop over ports. \code{plic.rs} writes
\code{Plic1} to \code{Plic8}; the macro takes counts from 1 to 31, the
most one word of pending bits holds. \code{plic.rs} also names the
offsets, \code{PENDING}, \code{ENABLE}, \code{THRESHOLD} and
\code{CLAIM}, and \code{priority(i)}.

\subsection*{Parameters}

\begin{center}\footnotesize
\begin{tabular}{@{}lp{0.7\textwidth}@{}}
\toprule
Parameter & Meaning \\
\midrule
\textit{N} & The count of sources, in the type's name. \\
\code{EDGE} & A bit per source: bit \textit{i} set makes source
\textit{i} ask on a rising edge rather than while its line is high.
Bit 0 and bits above \textit{N} are ignored. The tables use 0, as the
board does. \\
\bottomrule
\end{tabular}
\end{center}

\subsection*{Register map}

Offsets from the controller's base, which are the RISC-V PLIC's for
one target. The controller decodes the low 22 bits of an address. On
the board the base is \code{0x0c00\_0000}.

\begin{center}\footnotesize
\begin{tabular}{@{}llp{0.55\textwidth}@{}}
\toprule
Offset & Word & Access \\
\midrule
\code{0x4} $\times$ \textit{i} & Source \textit{i}'s priority, 3 bits &
Read and write, for \textit{i} from 1 to \textit{N}. \\
\code{0x1000} & The pending bits, bit \textit{i} for source \textit{i} &
Read only. \\
\code{0x2000} & The enable bits, the same way & Read and write; bit 0
does not stick. \\
\code{0x20\_0000} & The threshold, 3 bits & Read and write. \\
\code{0x20\_0004} & Claim and complete & A read claims; a write
completes. \\
\bottomrule
\end{tabular}
\end{center}

Every other offset, offset 0 among them, reads as zero and ignores a
write.

\dstables{Plic}

\subsection*{Behaviour}

\begin{itemize}
\item \textbf{The gateway.} A level source asks while its line is high.
An edge source asks when its line is high and was low a cycle before
(\code{prev}), or when it holds an edge (\code{held}).
\item A request goes forward when its source is not \code{active}.
Going forward sets the source's bits in \code{pending} and in
\code{active}.
\item An edge that comes while its source is \code{active} sets its bit
in \code{held}, one deep, and it asks as soon as the source is no
longer active.
\item \textbf{The choice.} A source is a candidate when it is pending,
enabled, and its priority is above \code{threshold}. The best
candidate has the highest priority, and the lowest number wins a tie.
A priority of 0 is never above a threshold, so it never interrupts.
\item \textbf{A claim.} A read of \code{0x20\_0004} answers with the
best candidate's number, or 0 for none, and clears that source's
\code{pending} bit. The source stays \code{active}.
\item \textbf{A complete.} A write of a number from 0 to \textit{N} to
\code{0x20\_0004} clears that source's \code{active} bit. A larger
number changes nothing. A level source whose line is still high then
asks again at once.
\item \textbf{The line.} \code{asserted} is set each cycle when there
is a best candidate, and \code{irq} follows \code{asserted}, so the
line is a register's output and lags the choice by a cycle.
\item \textbf{The bus.} A read is answered in the cycle it is taken, and
a write is taken when its address and its word are both there, and
answered \code{Okay} at once. Every answer is \code{Okay}.
\end{itemize}

\subsection*{Verification}

\begin{itemize}
\item In \code{plic}, in \code{//lib/parts:parts\_test}, against a
\code{Plic3}:
\begin{itemize}
\item \code{the\_highest\_priority\_is\_claimed\_first\_and\_the\_lowest\_number\_wins\_a\_tie};
\item \code{a\_priority\_at\_or\_below\_the\_threshold\_does\_not\_interrupt};
\item \code{a\_claimed\_source\_asks\_again\_only\_after\_its\_complete};
\item \code{an\_edge\_source\_asks\_once\_per\_edge\_and\_a\_level\_source\_while\_high}.
\end{itemize}
\code{the\_eight\_source\_controller\_lowers} writes \code{Plic8}'s
Verilog.
\item \code{ex\_plic} runs \code{Plic3<0b100>} against a host that
claims and completes, and the build simulates its netlist against that
run under nvc and Verilator: \code{//docs:plic\_sim\_plic3\_tb\_test}
and \code{//docs:plic\_vsim\_test}.
\item \code{input\_comes\_by\_interrupt\_one\_byte\_each}, in
\code{//cpu/vreteno:board\_test}, runs the board's \code{Plic2} under a
compiled program that takes the serial port's input by interrupt.
\end{itemize}

\subsection*{Used by}

\code{Plic2<0>} in the board's design, \code{Board}, as its field
\code{plic} behind the bridge \code{pplic}: source 1 is the serial
port's receive interrupt, source 2 the board's \code{irq} input, and
\code{irq} is the core's external interrupt. \code{Plic3<0b100>} in the
example \code{ex\_plic}.

\subsection*{Limits}

It has one target and no per-target context beyond it. Priorities are 3
bits, so from 0 to 7. A write ignores its strobe and writes the whole
field. The controller checks no address above bit 21, and relies on its
bridge for the rest. An edge source holds one edge during a service; a
second is lost.

% SPDX-License-Identifier: Apache-2.0
% covers: Gpio
\datasheet{Gpio}{General purpose pins on AXI-Lite: a value out, a value in, a direction, and an interrupt per pin.}

\dsfacts{\code{lib/parts/src/gpio.rs}}{\code{txhdl\_parts::gpio::Gpio}}%
{\code{//docs:examples}}

\subsection*{Function}

\code{Gpio<N>} puts \code{N} pins behind an AXI-Lite link. A program
chooses which way each pin faces, drives the ones that face out, reads
the ones that face in, and may ask for an interrupt from any of them.

The part drives no pin itself. It offers the three wires a pad wants,
\code{drive}, the value; \code{dirs}, whether to drive it; and
\code{pins}, the value read back, and a board's top level joins them
to a tristate buffer. A \code{Pad} carries nothing in a simulation, so
a part written around one could not be run or checked.

\code{pins} is taken through two flip-flops, \code{sync0} then
\code{sync1}, before anything reads it. Nothing else in the part reads
the port.

\subsection*{Parameters}

\begin{center}\footnotesize
\begin{tabular}{@{}lp{0.7\textwidth}@{}}
\toprule
Parameter & Meaning \\
\midrule
\code{N} & How many pins there are. Every register is \code{N} bits
wide and reads back into the low \code{N} bits of a word; the bits
above are zero. The tables use 8, as \code{ex\_gpio} does. \\
\bottomrule
\end{tabular}
\end{center}

\subsection*{Register map}

Offsets from the peripheral's base. The part selects the word by bits
4 to 2 of the address and looks at no other bits, so the router or the
bridge in front of it decides the base.

\begin{center}\footnotesize
\begin{tabular}{@{}llp{0.5\textwidth}@{}}
\toprule
Offset & Word & Access \\
\midrule
\code{0x00} & \code{out} & Read and write. Driven on a pin whose
direction is out. \\
\code{0x04} & \code{sync1} & Read only. What the pins read, after the
two flip-flops. A write is taken and ignored. \\
\code{0x08} & \code{dir} & Read and write. One to drive the pin, zero
to read it. \\
\code{0x0c} & \code{ie} & Read and write. One to let the pin raise
\code{irq}. \\
\code{0x10} & \code{kind} & Read and write. Zero for a level, one for
an edge. \\
\code{0x14} & \code{pol} & Read and write. A level: one for high. An
edge: one for rising. \\
\code{0x18} & \code{status} & Read, and write one to clear. Which pins
have fired. \\
\bottomrule
\end{tabular}
\end{center}

VHDL reserves \code{out}, so the netlist calls the register
\code{out\_rw} in both targets (issue 497).

\dstables{Gpio}

\subsection*{Behaviour}

\begin{itemize}
\item Every cycle \code{sync0} takes \code{pins}, \code{sync1} takes
\code{sync0}, and \code{seen} takes \code{sync1}. So a read of
\code{sync1} is what the pin held two cycles before, and \code{seen}
is what it held three cycles before, which is what an edge is measured
against.
\item A pin fires while it is enabled in \code{ie} and either its
\code{kind} bit is zero and \code{sync1} equals \code{pol}, or its
\code{kind} bit is one and \code{sync1} crossed \code{seen} in the
direction \code{pol} names.
\item \code{status} takes what it held, less any bit a write to
\code{0x18} set to one in the same cycle, plus every pin firing that
cycle. A level therefore sets its bit again in the cycle after it is
cleared, while the level lasts. Clearing is by a written one and not a
written zero, so two programs clearing different pins cannot lose each
other's work.
\item \code{irq} is high while any bit of \code{status} is set. It is
combinational from \code{status}, which is a register, so it follows
the register by no cycles.
\item A read is answered on \code{bus\_r} in the cycle it is taken, with
\code{OKAY}. A write is taken when its address phase, its beat and
room on \code{bus\_b} are all there, and is answered \code{OKAY} in that
cycle.
\item An address that selects no word reads as zero and a write to it
is answered and does nothing.
\end{itemize}

\subsection*{Verification}

\begin{itemize}
\item \code{//lib/examples:ex\_gpio} sets the direction, drives a
pattern, reads a pin through the synchroniser, takes an interrupt on a
rising edge and not on the fall, clears the status, and asks for a
level and finds the bit back. Every step is asserted against what the
program read.
\item The build simulates \code{Gpio<8>}'s netlist against that run
under nvc and Verilator: \code{//docs:gpio\_sim\_gpio\_tb\_test} and
\code{//docs:gpio\_vsim\_test}.
\end{itemize}

\subsection*{Used by}

Nothing in the tree yet. It is meant for a board's LEDs and keys, and
for the header pins, behind an AXI-Lite bridge; issue 145 says so.

\subsection*{Limits}

The interrupt is one line for all \code{N} pins; which pin raised it
is in \code{status}. There is no debounce, so a key wired to a pin
bounces into an edge interrupt many times. The part does not drive a
pad, so a board that wants one writes the tristate buffer itself.
\code{N} is at most the width of a bus word, since every register
reads back in one.

% SPDX-License-Identifier: Apache-2.0
% covers: Spi
\datasheet{Spi}{An SPI master on AXI-Lite: one byte each way at a time, under a chip select a program holds.}

\dsfacts{\code{lib/parts/src/spi.rs}}{\code{txhdl\_parts::spi::Spi}}%
{\code{//docs:examples}}

\subsection*{Function}

\code{Spi} drives the four wires of an SPI link: \code{sclk}, the
clock; \code{mosi}, the byte going out; \code{miso}, the byte coming
in; and \code{cs\_n}, the chip select, low while a chip is held.

A program writes a byte to \code{data}, waits for the transfer to
finish, and reads \code{data} for the byte that arrived the other way.
SPI has no direction: the eight clocks that send a byte are the eight
that receive one, so every transfer is a swap, and a driver that wants
to read sends a byte it does not care about.

The chip select is a bit in \code{ctrl} that a program sets and
clears, not something a transfer does. Every SPI command is several
bytes under one unbroken select, so a master that dropped it between
bytes would start a new command each time.

\subsection*{Parameters}

None. The divider and the mode are registers rather than parameters,
so one instance can reach chips that want different clocks.

\subsection*{Register map}

Offsets from the peripheral's base. The part selects the word by bits
3 and 2 of the address and looks at no other bits.

\begin{center}\footnotesize
\begin{tabular}{@{}llp{0.5\textwidth}@{}}
\toprule
Offset & Word & Access \\
\midrule
\code{0x00} & \code{ctrl} & Read and write. Bits 7 to 0 the divider,
8 \code{cpol}, 9 \code{cpha}, 10 the chip select, 11 the interrupt
enable. \\
\code{0x04} & \code{data} & Written: the byte to send, which starts a
transfer. Read: the byte that came in. \\
\code{0x08} & \code{state} & Bit 0 a transfer is running, bit 1 one
has finished. Write a one to bit 1 to clear it. \\
\bottomrule
\end{tabular}
\end{center}

A half of a bit takes \code{div + 1} cycles, so the clock is the
system's over \code{2 * (div + 1)}.

\dstables{Spi}

\subsection*{Behaviour}

\begin{itemize}
\item A write to \code{data} while no transfer is running loads
\code{txb}, sets \code{busy}, and clears \code{tick}, \code{half} and
\code{count}. A write while one is running is answered and does
nothing, so a driver reads \code{state} first.
\item While \code{busy}, \code{tick} counts to \code{div}. The cycle
it reaches it, \code{half} turns over and \code{count} adds one.
\code{sclk} is \code{cpol} exclusive-or \code{half}, so it idles at
\code{cpol}.
\item The bit is taken on the leading turn when \code{cpha} is clear
and on the trailing turn when it is set. \code{rxb} shifts left and
takes \code{miso}.
\item \code{txb} shifts on the other turn, with one exception: when
\code{cpha} is set the first leading turn does not shift, because the
first bit is already on the wire and moving it would lose the last bit
of the byte. That is the turn at \code{count} zero.
\item After sixteen turns, eight bits, \code{busy} clears,
\code{half} returns to zero and \code{fired} is set. \code{irq} is
\code{fired} and the interrupt enable.
\item \code{mosi} is the top bit of \code{txb} and \code{cs\_n} is the
inverse of the select bit, both combinational.
\item A read is answered on \code{bus\_r} in the cycle it is taken, with
\code{OKAY}. A write is taken when its address phase, its beat and
room on \code{bus\_b} are all there.
\end{itemize}

\subsection*{Verification}

\begin{itemize}
\item \code{//lib/examples:ex\_spi} reads the chip's identity with
\code{0x9f} and four bytes from an address with the \code{0x0b} fast
read, against \code{FlashDevice}, and asserts both against what was
put in the model.
\item The build simulates \code{Spi}'s netlist against that run under
nvc and Verilator: \code{//docs:spi\_sim\_spi\_tb\_test} and
\code{//docs:spi\_vsim\_test}.
\end{itemize}

\subsection*{Used by}

Nothing in the tree yet. It is meant for the board's configuration
flash, for an SD card in SPI mode, and for sensors on the header;
issue 146 says so.

\subsection*{Limits}

One chip select, so a second chip wants a second instance or a pin
from elsewhere. One byte at a time, so a driver spends a bus
transaction per byte and the interrupt is per byte, not per command.
No quad or dual mode, and no direct memory access: reading a large
image through this costs the core a read and a write for every byte.
The board's configuration flash is reached through the
\code{STARTUPE2} primitive rather than an ordinary pin, which this
part does not do.

% SPDX-License-Identifier: Apache-2.0
% covers: Syscon
\datasheet{Syscon}{The system control block: identity, build stamp, reset cause, soft reset, and a word that survives a reset.}

\dsfacts{\code{lib/parts/src/syscon.rs}}{\code{txhdl\_parts::syscon::Syscon}}%
{\code{//docs:examples}}

\subsection*{Function}

\code{Syscon} answers three questions only the hardware can answer:
which design is running, why it last restarted, and where to leave a
word for whatever runs next. It also drives the system's reset line.

Four of its words are constants from the type's parameters, so a build
can put a commit and a date into the hardware and a program can read
back which bitstream it is on. Three are registers: the cause, the
reset key and the scratch.

\subsection*{Parameters}

\begin{center}\footnotesize
\begin{tabular}{@{}lp{0.68\textwidth}@{}}
\toprule
Parameter & Meaning \\
\midrule
\code{ID} & The design's identifier, read at \code{0x00}. \\
\code{VERSION} & Its version, read at \code{0x04}. \\
\code{STAMP0} & The build stamp's low word, read at \code{0x08}. A
build would put a commit here. \\
\code{STAMP1} & The build stamp's high word, read at \code{0x0c}. A
build would put a date here. \\
\code{KEY} & The word that must be written to \code{0x14} for a reset
to happen. Any other value is ignored. \\
\bottomrule
\end{tabular}
\end{center}

\subsection*{Register map}

Offsets from the block's base. The block selects the word by bits 4 to
2 of the address and looks at no other bits.

\begin{center}\footnotesize
\begin{tabular}{@{}llp{0.5\textwidth}@{}}
\toprule
Offset & Word & Access \\
\midrule
\code{0x00} & \code{id} & Read only, \code{ID}. \\
\code{0x04} & \code{version} & Read only, \code{VERSION}. \\
\code{0x08} & \code{stamp0} & Read only, \code{STAMP0}. \\
\code{0x0c} & \code{stamp1} & Read only, \code{STAMP1}. \\
\code{0x10} & \code{cause} & Read, and write ones to clear. Bit 0 the
power, 1 the button, 2 the watchdog, 3 software. \\
\code{0x14} & \code{reset} & Write only. Writing \code{KEY} asks for a
reset; anything else does nothing. Reads as zero. \\
\code{0x18} & \code{scratch} & Read and write. A reset does not clear
it. \\
\bottomrule
\end{tabular}
\end{center}

\dstables{Syscon}

\subsection*{Behaviour}

\begin{itemize}
\item Each cycle \code{cause} takes what it held, less any bit a write
to \code{0x10} set to one in that cycle, plus a bit for each reason
asserting now: \code{por}, \code{button}, \code{wdog}, and a write of
\code{KEY} to \code{0x14}.
\item The reasons accumulate. A system reset by the button and then by
software reads both bits until a program clears them, so the case
where two things went wrong is not lost.
\item A write of \code{KEY} to \code{0x14} sets \code{hold} to eight.
While \code{hold} is not zero it counts down by one a cycle. The reset
is held for eight cycles rather than one so that every part of a
design sees it.
\item \code{srst} is high while any of \code{por}, \code{button},
\code{wdog} or the key write is asserting, or while \code{hold} is not
zero. It is combinational from those and from a register.
\item Nothing in this block is cleared by \code{srst}. That is what
lets \code{cause} and \code{scratch} survive a reset, and it is why
the block sits outside whatever \code{srst} drives.
\item A read is answered on \code{bus\_r} in the cycle it is taken, with
\code{OKAY}. A write is taken when its address phase, its beat and
room on \code{bus\_b} are all there.
\end{itemize}

\subsection*{Verification}

\begin{itemize}
\item \code{//lib/examples:ex\_syscon} reads the identity and the
build stamp, reads the power-on cause, leaves a word in
\code{scratch}, clears the cause, asks for a reset with the key, and
checks that the cause comes back as software while the scratch word is
unchanged. It then writes a wrong key and checks that nothing moved,
and presses the button to see the two causes side by side.
\item The build simulates the netlist against that run under nvc and
Verilator: \code{//docs:syscon\_sim\_syscon\_tb\_test} and
\code{//docs:syscon\_vsim\_test}.
\end{itemize}

\subsection*{Used by}

Nothing in the tree yet. Issue 147 asks for it on the board, where the
watchdog of issue 149 would drive \code{wdog} and the bootloader of
issue 135 would use \code{scratch}.

\subsection*{Limits}

The four constants are type parameters, so putting a real commit in
them wants a build that generates the instantiation; this part offers
the words and does not fill them. The reset is eight cycles, not a
parameter. There is no lock, so a program can clear a cause another
program wanted to read. \code{reset} reads as zero rather than as what
was last written.

% SPDX-License-Identifier: Apache-2.0
% covers: Uart
\datasheet{Uart}{Vreteno's serial port: an AXI-Lite peripheral with a line each way and a receive buffer of eight.}

\dsfacts{\code{cpu/vreteno/src/uart.rs}}{\code{vreteno32::uart::Uart}}%
{\code{//docs:vreteno}, \code{//docs:noc}}

\subsection*{Function}

\code{Uart} sends bytes on \code{tx} and receives them on \code{rx},
one start bit, eight data bits least significant first and one stop
bit, each \code{DIV} cycles long. The lines rest high. A received byte
lands in a buffer of eight, and \code{irq} is high while the buffer
holds any.

It is an AXI-Lite peripheral. It sits behind a \code{LiteBridge1},
which hands it one transaction at a time with no identifier and no
burst. Reads come on \code{ar} and are answered on \code{r}; writes
come on \code{aw} and \code{w} and are answered on \code{b}. The input
\code{rst} stops both sides and empties the buffer.

\subsection*{Parameters}

\begin{center}\footnotesize
\begin{tabular}{@{}lp{0.7\textwidth}@{}}
\toprule
Parameter & Meaning \\
\midrule
\code{DIV} & Cycles per bit on both lines. The tables use 868, which
\code{uart.rs} states is 115200 baud at 100~MHz, as the board runs.
The checked runs use 4, and \code{board\_netlist sim} uses 16. \\
\bottomrule
\end{tabular}
\end{center}

\subsection*{Register map}

Offsets from the port's base. On the board the base is \code{0x3000},
which is also \code{isa::UART\_BASE}. The port selects the word by
bits 3 and 2 of the address.

\begin{center}\footnotesize
\begin{tabular}{@{}lp{0.38\textwidth}p{0.38\textwidth}@{}}
\toprule
Offset & Read & Write \\
\midrule
\code{0x0} & The last byte accepted for sending, zero-extended. &
Bits 7 to 0 start a frame if the port is idle; else the byte is
dropped. \\
\code{0x4} & Status: bit 0 busy sending, bit 1 a byte received and not
yet read, bit 2 the buffer full. Other bits zero. & Ignored. \\
\code{0x8} & The oldest received byte, zero-extended. The read removes
it from the buffer. & Ignored. \\
\code{0xc} & Zero. & Ignored. \\
\bottomrule
\end{tabular}
\end{center}

\dstables{Uart}

\subsection*{Behaviour}

\begin{itemize}
\item A read is taken when \code{r} has room and a request is on
\code{ar}, and it is answered with \code{OKAY} in the same cycle.
\item A write is taken when \code{b} has room and both \code{aw} and
\code{w} hold a request. It is answered with \code{OKAY} in the same
cycle, whichever word it names.
\item A byte written to word 0 while the port is idle loads the ten
bits of its frame into \code{shift} and ten into \code{bits}. It also
goes into \code{last}, and \code{sent} counts it.
\item While \code{bits} is not zero the port is busy and \code{tx} is
bit 0 of \code{shift}. Every \code{DIV} cycles the frame shifts down
and \code{bits} counts down. When idle, \code{tx} is high.
\item The \code{rx} line passes through one register, \code{line},
before the logic.
\item When the port is not receiving and \code{line} is low, a frame
begins. The port samples \code{line} in cycle \code{DIV/2} of each of
the ten bits.
\item A high where the start bit should be is a false start, and the
frame is dropped. The eight data bits shift into \code{rx\_shift}.
\item A high stop bit lands the byte. If the buffer holds eight, the
byte is dropped and \code{dropped} counts it; else it goes in at the
tail and \code{received} counts it. A low stop bit lands nothing.
\item A read of word 2 with a byte in the buffer removes the oldest. A
push and a removal in the same cycle leave the count unchanged.
\item \code{irq} is high while the buffer's count is not zero.
\item With \code{rst} high, \code{bits}, \code{tick},
\code{rx\_bits}, \code{rx\_tick}, \code{head} and \code{count} go to
zero.
\end{itemize}

\subsection*{Verification}

\begin{itemize}
\item \code{//cpu/vreteno:lockstep} puts the terminal of
\code{term.rs} on the two lines. At the halt it checks \code{sent} and
\code{last} against the bytes the model was given, and that
\code{dropped} is zero. In \code{demo\_program} the port must say
\code{OK} and a newline, receive three bytes, and echo them.
\item \code{//cpu/vreteno:hello\_test},
\code{//cpu/vreteno:hello\_cc\_test} and \code{//cpu/vreteno:board\_test}
check what compiled programs print on \code{tx}.
\item The build simulates \code{Uart<4>}'s netlist, the entity
\code{uart4}, against the trace of \code{//cpu/vreteno:vreteno} under
nvc and Verilator: \code{//docs:vreteno\_sim\_uart4\_uart4\_tb\_test}
and \code{//docs:vreteno\_vsim\_uart4\_test}. The run also writes
\code{Uart<868>}'s netlist, which no test simulates.
\item \code{//docs:vreteno} reports that the board, programmed with the
design from before the DDR3 memory was on the bus, said \code{OK} at
115200 baud and echoed the three bytes a watcher typed.
\end{itemize}

\subsection*{Used by}

\code{Board}, as \code{uart} behind the bridge \code{puart}. The
demonstration \code{//cpu/vreteno:vreteno}, the lockstep test and
\code{run.rs}. The system in \code{soc}, where it is the serial port at
corner 1,1 of the lattice.

\subsection*{Limits}

A byte written while a frame is going out is dropped, and a byte
received with the buffer full is dropped and counted. A program polls
the status and keeps up. \code{sent}, \code{received} and
\code{dropped} are eight-bit counters that are not on the bus. The port
checks no address beyond bits 3 and 2, since the bridge sends it only
its own range. The frame is fixed at ten bits, with no parity bit, and
the port has no flow-control lines.

% SPDX-License-Identifier: Apache-2.0
% covers: Ddr3Per
\datasheet{Ddr3Per}{The board's DDR3 memory as an AXI peripheral: the Wishbone bridge and AMD's MIG controller, which makes the design's clock.}

\dsfacts{\code{ddr3/src/lib.rs}}{\code{ddr3::Ddr3Per}}%
{\code{//docs:vreteno}, \code{//docs:axi}}

\subsection*{Function}

\code{Ddr3Per} puts the DDR3 memory of the Alinx AX7A200 on an AXI
link. It is a unit of units with two children. \code{bridge} is
\code{AxiWb} of \code{txhdl\_parts}, which serves a burst a word at a
time as requests on a pipelined Wishbone bus. \code{ctl} is
\code{Ddr3}, AMD's MIG 7 Series controller, which the build generates
as \code{//ddr3:ddr3\_mig}, behind the wrapper
\code{ddr3/hdl/ddr3\_wb32.v}, which gives it a Wishbone of 32-bit
words. Nine Wishbone signals join the two.

\code{Ddr3} is a foreign unit. Its netlist is an instance of the
module \code{ddr3\_wb32}, which the lowering names and does not write.
Its \code{run} is a model for simulation: a Wishbone memory that
stalls while it calibrates. The model does not drive the memory's
pins.

On the link side the peripheral is the client of an \code{AxiPer},
and the link is one port, \code{bus}, a \code{PerPort}:
\code{bus\_req} and \code{bus\_w} in, \code{bus\_ans} and
\code{bus\_r} out. On the
memory side its ports are the board's clock and the controller's
reset in, the design's clock and its reset out, since the controller
makes them, the \code{calib} output, and the memory's pins, with
\code{dq}, \code{dqs} and \code{dqs\_n} as pads both ways.

\subsection*{Parameters}

None. The controller calibrates in hardware, and its simulation
library carries the variant with the fast calibration, so neither
wants a parameter. The bridge inside is
\code{AxiWb<32, 2, 28>}, with \code{AW} equal to 28.

\subsection*{Address map}

The bridge's Wishbone address is a word address of \code{AW}, 28
bits: the burst's byte address shifted down by two and cut to 28 bits.
\code{lib.rs} states the 28 bits as fifteen row bits, ten column bits
and three bank bits. The wrapper takes the controller's address above
three lane bits, and the lane selects one word of the controller's
burst of eight.

On the board, \code{BoardRouter} sends the peripheral every address
that matches \code{0x4000\_0000} under the mask \code{0xc000\_0000}.
The cut to 28 bits removes the base, so the memory sees offsets from
\code{0x4000\_0000}.

\dstables{Ddr3Per}

\subsection*{Behaviour}

\begin{itemize}
\item The bridge takes one burst at a time and serves it a word at a
time at consecutive word addresses. A read's words go on the lines
one after another and its beats go back as they are acknowledged,
the last of them marked; a write goes a beat at a time as its beats
arrive and is answered once after the last.
\item The bridge offers a request with \code{cyc} and \code{stb}, and
the controller takes it in a cycle with no \code{stall}. The bridge
then holds \code{cyc} alone until \code{ack}, which brings the word
for a read. The answer goes on the link when the link has room.
\item Every Wishbone line the bridge drives is a register.
\item The wrapper repeats the word across the controller's burst of
eight with the mask clear on that one word alone, so a write touches
only its word; a read keeps its lane until the burst comes back and
picks its word out of it. One request is in flight at a time, and a
write is acknowledged the cycle after the controller takes it.
\item The controller makes the clocks. It takes the board's 200~MHz on
\code{sys\_clk}, which is also its delay reference, runs the memory
at 400~MHz, and hands out \code{ui\_clk}, the 100~MHz the design
runs on, with \code{ui\_rst} high until that clock is good. In the
netlist the controller's \code{i\_ui\_clk} is the design's clock,
which is that same clock fed back. \code{sys\_rst} is active high, a
pulse at power-on on the board.
\item \code{calib} is the controller's \code{o\_calib\_complete}.
\item In simulation the model stalls for \code{MODEL\_WARMUP}, 64
cycles, and then answers each request a cycle after it takes it.
\code{calib} rises when the warm-up has passed.
\end{itemize}

\subsection*{Verification}

\begin{itemize}
\item \code{//ddr3:ddr3\_test} runs \code{words\_go\_in\_and\_come\_back}.
Behind an \code{AxiHost} and an \code{AxiPer}, it writes words at
\code{0x4000\_0000}, \code{0x4000\_0104} and \code{0x47ff\_fffc}, the
first while the model still calibrates. It reads them back and checks
that \code{calib} did not rise before the warm-up.
\item \code{//ddr3:ddr3\_test} also runs
\code{the\_controller\_is\_instantiated\_and\_not\_written}. It checks
that the Verilog and the VHDL of \code{Ddr3Per} instantiate
\code{ddr3\_wb32}. The controller's clock must be on \code{clk}, the
board's clock and the clock it makes on the ports, and its pads on
the ports. The bridge's module must be
written and the controller's must not.
\item \code{//cpu/vreteno:board\_test} runs
\code{the\_memory\_test\_runs\_on\_the\_board}, a compiled program that
writes and reads the memory through the model and must print
\code{ddr3 ok}, and \code{the\_netlist\_holds\_the\_controller}.
\item \code{//ddr3/sim:wrapper\_test} simulates the wrapper and the
controller against the Micron model under Vivado's simulator:
calibrated, two words written at different lanes and read back. It
is manual. \code{//ddr3/sim:example\_test} is AMD's own proof of the
controller as configured, the example design's traffic generator
writing the memory and reading it back through it; manual too.
\item \code{//cpu/vreteno/board/sim:}\allowbreak\code{board\_test} simulates the whole
board, this controller included, against two Micron models. It is
manual.
\end{itemize}

\subsection*{Used by}

\code{Board}, as \code{ddr3} behind the tracker \code{pddr3}, at
\code{0x4000\_0000}.

\subsection*{Limits}

The Rust run of \code{Ddr3} is a model and not the controller. A
lowered netlist needs \code{//ddr3:hdl}, the wrapper, beside it, and
the controller's library, \code{//ddr3:ddr3\_mig}, which the build
generates under Vivado and which stays out of the tree, being AMD's.
The bridge serves one burst at a time. \code{//docs:vreteno} states
that the design with MIG as the controller has a bitstream that meets
timing, and that its run on the board is owed: the run it reports was
made with the controller before it.

% SPDX-License-Identifier: Apache-2.0
% covers: Board
\datasheet{Board}{The Vreteno machine for the Alinx AX7A200, as one lowered unit of units.}

\dsfacts{\code{cpu/vreteno/src/board.rs}}{\code{vreteno32::board::Board}}%
{\code{//docs:vreteno}}

\subsection*{Function}

\code{Board} is everything that computes on the board, written as a
unit of units so that \code{\#[lower]} makes its netlist as one module
holding the rest. It holds the core and its tracker, a router of five,
and the peripherals behind its five ports. The data memory and the
timer each sit behind an \code{AxiPer}. Four small peripherals share
the page at \code{0x3000} behind the AXI-Lite bridge
\code{LiteBridge4}, a sixteenth of the page each: the serial port at
\code{0x3000}, the pulse width modulator at \code{0x3100}, whatever a
board hangs on the slot at \code{0x3200}, and the remote peripheral at
\code{0x3300}. The DDR3 memory sits behind an \code{AxiPer} and
holds UberDDR3's controller as a foreign module. The interrupt
controller, \code{Plic2}, sits behind a second \code{LiteBridge1}.

\code{run.rs} wires the same parts for a simulation. \code{Board} is
that wiring as a unit, so a board top has to supply only the clocks,
the resets and the pins. The top for the AX7A200 is
\code{cpu/vreteno/board/vreteno\_board.v}.

\subsection*{Parameters}

\begin{center}\footnotesize
\begin{tabular}{@{}lp{0.7\textwidth}@{}}
\toprule
Parameter & Meaning \\
\midrule
\code{DIV} & The serial port's cycles per bit, as \code{Uart<DIV>}.
The tables use 868, the value \code{board\_netlist board} writes for
115200 baud at 100~MHz. \\
\code{MICRON\_SIM} & The DDR3 controller's \code{MICRON\_SIM}: one to
shorten its power-on waits for the Micron model, zero on the board. \\
\code{BIST} & The DDR3 controller's \code{BIST\_MODE}, its self test
after calibration. \\
\bottomrule
\end{tabular}
\end{center}

\code{board\_netlist} lowers \code{Board<868, 0, 0>} for the board and
\code{Board<16, 1, 0>} for simulation. \code{//cpu/vreteno:board\_test}
runs \code{Board<4, 1, 0>}.

\subsection*{Address map}

\code{BoardRouter} states the map in its type, \code{Router5<32, 32, 4,
2, ...>}. A peripheral gets an address when the address under the mask
equals the base.

\begin{center}\footnotesize
\begin{tabular}{@{}lllp{0.35\textwidth}@{}}
\toprule
Base & Mask & Child & Behind \\
\midrule
\code{0x1000} & \code{0xffff\_f000} & \code{dmem}, \code{Dmem<2>} &
\code{pdmem}, \code{AxiPer} \\
\code{0x0200\_0000} & \code{0xffff\_0000} & \code{timer},
\code{Timer<2>} & \code{ptimer}, \code{AxiPer} \\
\code{0x3000} & \code{0xffff\_f000} & \code{uart}, \code{Uart<DIV>} &
\code{puart}, \code{LiteBridge4} at \code{0x3000} \\
& \code{0xffff\_ff00} & \code{pwm}, \code{Pwm} &
the same bridge, at \code{0x3100} \\
& \code{0xffff\_ff00} & the board's, on its own clock &
the same bridge, at \code{0x3200} \\
& \code{0xffff\_ff00} & \code{remote},
\code{Remote<REMOTE\_WAIT>} & the same bridge, at \code{0x3300} \\
\code{0x4000\_0000} & \code{0xc000\_0000} & \code{ddr3},
\code{Ddr3Per} & \code{pddr3}, \code{AxiPer} \\
\code{0x0c00\_0000} & \code{0xfc00\_0000} & \code{plic},
\code{Plic2<0>} & \code{pplic}, \code{LiteBridge1} at
\code{0x0c00\_0000} \\
\bottomrule
\end{tabular}
\end{center}

The router answers a burst that matches no range with \code{DecErr}.
The core's instruction memory is not on the bus, and the core sends no
request for an address below \code{0x1000}.

\dstables{Board}

\subsection*{Behaviour}

\begin{itemize}
\item The core's tracker is \code{AxiHost<32, 32, 4, 2, 4>}: 32-bit
addresses and words, four lanes, two-bit identifiers, four of them.
\item \code{rst} goes to the core, the timer, the serial port and the
interrupt controller. The DDR3 controller has its own active-low
\code{ddr3\_rst\_n}.
\item The interrupt controller's source 1 is the serial port's
interrupt line, and its source 2 is \code{irq}; both ask while high.
Its line is the core's external interrupt. The timer's line goes to
the core's \code{tirq}.
\item Of the core's outputs, only \code{halt} is a port. \code{instr}
and the \code{Writeback} go nowhere.
\item The join runs the timer before the core, since the core reads
the timer's line in the same step, and the serial port before the
interrupt controller before the core, for the same reason. \code{Ddr3Per} runs its controller
before its bridge for the same reason.
\item \code{calib} and the memory's pins are \code{Ddr3Per}'s,
brought out as the unit's own ports.
\item The slot at \code{0x3200} does not reach a field. It leaves as
five channel ports, \code{vaw}, \code{var} and \code{vw} out and
\code{vb} and \code{vr} back, because what sits there runs on a clock
of its own and the crossing is the board top's business. A board with
nothing there ties them off, and a program that touched \code{0x3200}
on such a board would wait for ever.
\item The remote peripheral at \code{0x3300} is the other way about:
\code{remote} and \code{link} are fields, because both run on the
core's clock. What leaves the unit is the frame stream, \code{net\_tx}
out and \code{net\_rx} in, a byte and a last bit each, which is what
the MAC speaks. \code{REMOTE\_DEV} is 1 and \code{REMOTE\_WAIT} is
two million cycles, 20~ms at 100~MHz, after which an unanswered
transaction is \code{SLVERR} on the bus.
\item \code{board\_netlist} puts a program in the core's
\code{imem} and in the data memory's \code{lane0} to \code{lane3} of
the netlist, since nothing on the machine loads them at run time.
\end{itemize}

The top, \code{vreteno\_board.v}, does the following around the
lowered \code{board}:

\begin{itemize}
\item One PLL on the board's 200~MHz differential clock makes 100~MHz
for the design, 400~MHz for \code{ddr3\_clk}, the same 400~MHz a
quarter cycle late for \code{ddr3\_clk\_90}, and 200~MHz for
\code{ref\_clk}.
\item The memory's reset follows the button and the PLL's lock. The
core's reset also waits for \code{calib}. Each passes through two
flip-flops on the design's clock.
\item It ties \code{irq} low.
\item Its four LEDs, lit when driven low, show the core halted, a
heartbeat from bit 25 of a counter, the memory calibrated, and the PLL
locked.
\end{itemize}

\subsection*{Verification}

\begin{itemize}
\item \code{//cpu/vreteno:board\_test} runs \code{Board<4, 1, 0>} in
Rust. \code{the\_greeting\_runs\_on\_the\_board} must print
\code{hello from rust} and halt within 8000 cycles.
\code{the\_memory\_test\_runs\_on\_the\_board} must print
\code{ddr3 ok} and halt within 40000 cycles.
\code{input\_comes\_by\_interrupt\_one\_byte\_each} runs a program
that takes four typed bytes through the interrupt controller, one
interrupt each, and must print \code{got ping in 4 interrupts} and
halt within 20000 cycles.
\code{the\_netlist\_holds\_the\_controller} checks that the Verilog
has the top, the core and the bridge as modules, instantiates
\code{ddr3\_wb32} without writing it, and has the \code{dq} pads.
\item \code{//cpu/vreteno/board/sim:board\_test} simulates the top, the
netlist of \code{//cpu/vreteno:board\_sim\_v}, the controller and two
Micron models under Vivado's simulator. It is manual.
\code{//docs:vreteno} reports that in that run the core says
\code{ddr3 ok} on the serial port and halts.
\item \code{//cpu/vreteno:vreteno\_board\_synth} and
\code{//cpu/vreteno:vreteno\_board\_pnr} synthesise, place and route
the design with the board's pins. Both are manual.
\code{//docs:vreteno} reports that every timing constraint is met, with
0.432~ns worst setup slack and 0.031~ns worst hold slack. It reports
8100 LUTs, 6959 flip-flops, four block RAM tiles and four DSP blocks.
\item \code{//cpu/vreteno:vreteno\_board\_prog} and
\code{//cpu/vreteno:vreteno\_board\_flash} program the FPGA and the
QSPI flash. Both are manual.
\end{itemize}

No waveform target simulates \code{Board}'s netlist against a trace.

\subsection*{Used by}

\code{//cpu/vreteno:board\_netlist}, whose genrules
\code{//cpu/vreteno:board\_v} and \code{//cpu/vreteno:board\_sim\_v}
write the netlist, and \code{vreteno\_board.v}, which instantiates it
as \code{board}. The test \code{cpu/vreteno/tests/board.rs}.

\subsection*{Limits}

\code{//docs:vreteno} states that this design, with the DDR3 memory on
the bus, has not yet been run on the board. The board run it reports
was made with an earlier design without the memory. The program is
fixed in the netlist. The board top ties \code{irq} low, so on the
board the interrupt controller's second source never asks. The core is the only host.

% Razboj, the GPU.
% SPDX-License-Identifier: Apache-2.0
% covers: Fb
\datasheet{Fb}{Razboj's framebuffer: a memory of one word per pixel behind an AXI peripheral end.}

\dsfacts{\code{gpu/razboj/src/fb.rs}}{\code{razboj::fb::Fb}}%
{\code{//docs:razboj}, \code{//docs:axi}}

\subsection*{Function}

\code{Fb} is \code{N} words of memory, \code{px}, behind the
peripheral end of an AXI link. A pixel is one word. The rasteriser
writes pixels into it, and in Razboj's runs the same memory also
holds the display list and its count, which the rasteriser reads.

It is the peripheral client of an \code{AxiPer}: requests come on
\code{req} and write beats on \code{wd}, and answers go out on
\code{ans} and \code{rb}. It answers single-beat bursts, which is all
the rasteriser makes.

\subsection*{Parameters}

\begin{center}\footnotesize
\begin{tabular}{@{}lp{0.7\textwidth}@{}}
\toprule
Parameter & Meaning \\
\midrule
\code{A} & The width of the link's address, in bits. \\
\code{I} & The width of the AXI identifier, in bits. \\
\code{N} & Words of memory, a power of two. \\
\bottomrule
\end{tabular}
\end{center}

The tables use \code{Fb<16, 2, 1024>}, the framebuffer of
\code{//gpu/razboj:trace}.

\subsection*{Address map}

The word a burst names is its address shifted down by two, cut to
sixteen bits, and masked with \code{N}~$-$~1. The memory itself gives
no word a meaning. In \code{//gpu/razboj:trace} and the tests of
\code{gpu/razboj/src/sim.rs}, with \code{N} of 1024:

\begin{center}\footnotesize
\begin{tabular}{@{}lp{0.6\textwidth}@{}}
\toprule
Byte address & What \code{sim.rs} puts there \\
\midrule
\code{0x000} & The framebuffer, sixteen by sixteen pixels, a word
each. \\
\code{0x400} & The display list, eight words an instruction. \\
\code{0x600} & The count of instructions, written last. \\
\bottomrule
\end{tabular}
\end{center}

\dstables{Fb}

\subsection*{Behaviour}

\begin{itemize}
\item A read is taken when \code{rb} has room and no write is waiting
for its beat. It is answered on \code{rb} in the cycle it is taken,
with the word from \code{px}, \code{OKAY} and \code{last} set.
\item A write is taken when no write is waiting. Its word and its
identifier go into \code{paddr} and \code{pid}, and \code{pend} is
set.
\item The held write's beat is taken when it has arrived and
\code{ans} has room. The whole beat is written to the word, and
\code{ans} sends \code{OKAY} with \code{pid}. The strobe is not read.
\item While a write waits for its beat, no request of either kind is
taken.
\item An address beyond the memory wraps rather than ending the run.
\item \code{pixel} reads a word back for a run or a test.
\end{itemize}

\subsection*{Verification}

\begin{itemize}
\item \code{//gpu/razboj:razboj\_test} renders every scene on the rasteriser
through the link into \code{Fb}, reads the pixels back and compares
each with the model in \code{gpu/razboj/src/model.rs}. The tests are
\code{a\_scene\_of\_every\_kind\_agrees\_with\_the\_model},
\code{a\_clear\_says\_only\_its\_colour},
\code{a\_triangle\_is\_the\_same\_whichever\_way\_it\_is\_wound},
\code{a\_triangle\_hanging\_off\_the\_screen\_is\_clipped},
\code{a\_single\_pixel\_and\_an\_empty\_box} and
\code{pseudorandom\_scenes\_agree\_with\_the\_model}.
\item The build simulates \code{Fb<16, 2, 1024>}'s netlist, with the
display list and its count in \code{px}, against the trace of
\code{//gpu/razboj:trace} under nvc and Verilator:
\code{//docs:razboj\_sim\_fb\_fb\_tb\_test} and
\code{//docs:razboj\_vsim\_fb\_test}.
\item \code{//gpu/razboj:trace} and \code{//gpu/razboj:demo} assert that the
framebuffer equals the model's picture. The build runs
\code{//gpu/razboj:demo} to draw the picture in \code{//docs:razboj}.
\end{itemize}

\subsection*{Used by}

\code{razboj::sim::run} and \code{razboj::sim::run\_list}, and so Razboj's
tests, \code{//gpu/razboj:trace} and \code{//gpu/razboj:demo}. The system's
\code{//soc:demo} draws the core's display list a second time with
\code{run\_list}, on a plain link into \code{Fb}, and checks that the
picture is the same. On the lattice itself the memory is the
simulation-only \code{Ram}.

\subsection*{Limits}

It serves bursts of one beat and does not look at a burst's length or
strobe. It holds one write at a time. \code{//docs:razboj} states that
this is why the demonstration takes about three cycles a pixel. The
word index is sixteen bits wide, so the framebuffer reaches at most
65536 words.

% SPDX-License-Identifier: Apache-2.0
% covers: Raster
\datasheet{Raster}{Razboj's rasteriser: it reads a display list over AXI and writes one pixel a cycle.}

\dsfacts{\code{gpu/razboj/src/raster.rs}}{\code{razboj::raster::Raster}}%
{\code{//docs:razboj}, \code{//docs:noc}}

\subsection*{Function}

\code{Raster} finds its work in memory and draws it there. It polls a
count at \code{CTRL} until the count is not zero. It then reads the
six words of each instruction of the display list at \code{DL}, one
read at a time, and walks the instruction's box one pixel per cycle.
It writes the instruction's colour at every pixel inside the
primitive: every pixel of a clear or a rectangle, and every pixel of a
triangle at which no edge function is negative.

It is an AXI host written as hardware. It issues bursts on
\code{issue}, sends a pixel's beat on \code{wbeat}, and hands
identifiers back on \code{release}. It receives \code{grant},
\code{done} and \code{rdata} from its tracker. \code{idle} goes high
when the list is drawn and no write is in flight.

\subsection*{Parameters}

\begin{center}\footnotesize
\begin{tabular}{@{}lp{0.7\textwidth}@{}}
\toprule
Parameter & Meaning \\
\midrule
\code{A} & The width of the link's address, in bits. \\
\code{I} & The width of the AXI identifier, in bits. \\
\code{LOGW} & The screen is \code{1 << LOGW} pixels wide. \\
\code{H} & The screen is \code{H} pixels high. \\
\code{BASE} & The byte address of the framebuffer's first word. \\
\code{DL} & The byte address of the display list's first word. \\
\code{CTRL} & The byte address of the count. \\
\bottomrule
\end{tabular}
\end{center}

The tables use \code{Raster<16, 2, 4, 16, 0, 0x400, 0x600>}, the
rasteriser of \code{//gpu/razboj:trace}: a sixteen by sixteen screen at
\code{0x0}, the list at \code{0x400} and the count at \code{0x600}. In
\code{soc} it is \code{Raster<32, 2, 7, 96, 0x8000, 0x2000, 0x4000>}.

\subsection*{Address map}

\begin{center}\footnotesize
\begin{tabular}{@{}lp{0.6\textwidth}@{}}
\toprule
Address & What the rasteriser does there \\
\midrule
\code{CTRL} & Reads the count while polling. \\
\code{DL + (i << 5) + (k << 2)} & Reads word \code{k}, 0 to 5, of
instruction \code{i}. \\
\code{BASE + (((y << LOGW) + x) << 2)} & Writes pixel \code{x},
\code{y}. \\
\bottomrule
\end{tabular}
\end{center}

An instruction's words are in the format \code{gpu/razboj/src/dl.rs} states.
Words 6 and 7 of an instruction are never read.

\begin{center}\footnotesize
\begin{tabular}{@{}lp{0.42\textwidth}p{0.3\textwidth}@{}}
\toprule
Word & Low field & High field \\
\midrule
0 & bits 1 to 0: kind, 0 clear, 1 rectangle, 2 triangle & bits 25 to
2: colour \\
1 & bits 9 to 0: \code{x0} & bits 25 to 16: \code{y0} \\
2 & bits 9 to 0: \code{x1} & bits 25 to 16: \code{y1} \\
3 & bits 11 to 0: \code{ax} & bits 27 to 16: \code{ay} \\
4 & bits 11 to 0: \code{bx} & bits 27 to 16: \code{by} \\
5 & bits 11 to 0: \code{cx} & bits 27 to 16: \code{cy} \\
\bottomrule
\end{tabular}
\end{center}

\dstables{Raster}

\subsection*{Behaviour}

\begin{itemize}
\item \code{phase} is 0 while polling, 1 while fetching, 2 while
drawing and 3 when done.
\item While polling or fetching, the rasteriser issues one read when
no read is out and \code{issue} has room. \code{waiting} is set until
the read's beat arrives.
\item A count that is not zero starts the fetch at instruction 0, word
0, and its low eight bits go into \code{left}. A count of zero means
poll again.
\item During the fetch each word's fields are latched as the word
lands. A kind field of 3 is taken as a triangle.
\item When word 5 lands, the walk is set up in that cycle. The third
vertex comes straight from the bus. The six products of the edge
functions are computed then, and nowhere else.
\item A clear's box is the whole screen, from its own type. A
rectangle's and a triangle's box is \code{x0}, \code{y0} to
\code{x1}, \code{y1}. Vertices are sign-extended from twelve bits.
\item Each edge is kept as its value at this pixel, its value at the
start of this row and its two steps, in 32-bit two's complement. A
step to the right adds the column step; a new row reloads from the row
value plus the row step.
\item \code{hit} is high when the kind is not a triangle, or when all
three edge values have bit 31 clear.
\item A pixel with \code{hit} high is written when both \code{issue}
and \code{wbeat} have room. The walk steps when the pixel needs no
write or its write went out, so backpressure holds the walk.
\item A pixel's write is a burst of one beat, size 2, \code{INCR},
with the colour zero-extended to 32 bits, strobe \code{0xf} and
\code{last} set.
\item At the end of the box the rasteriser fetches the next instruction
while \code{left} is above 1, and otherwise goes to phase 3.
\item \code{inflight} counts writes issued whose response has not come
back.
\item When a write response and a read beat arrive together, the
write's identifier goes back first and the read's waits a cycle.
Grants are taken and dropped.
\end{itemize}

\subsection*{Verification}

\begin{itemize}
\item \code{//gpu/razboj:razboj\_test} renders scenes on the rasteriser through
an \code{AxiHost}, an \code{AxiPer} and \code{Fb}, and compares every
pixel with the model of \code{gpu/razboj/src/model.rs}:
\code{a\_scene\_of\_every\_kind\_agrees\_with\_the\_model},
\code{a\_clear\_says\_only\_its\_colour},
\code{a\_triangle\_is\_the\_same\_whichever\_way\_it\_is\_wound},
\code{a\_triangle\_hanging\_off\_the\_screen\_is\_clipped},
\code{a\_single\_pixel\_and\_an\_empty\_box}, and
\code{pseudorandom\_scenes\_agree\_with\_the\_model}, twelve scenes of
up to four random rectangles and triangles.
\item \code{an\_instruction\_survives\_the\_format} in the same target
checks the display list's encoder against its decoder.
\item The build simulates the netlist of \code{Raster<16, 2, 4, 16, 0,
0x400, 0x600>} against the trace of \code{//gpu/razboj:trace} under nvc and
Verilator: \code{//docs:razboj\_sim\_raster\_raster\_tb\_test} and
\code{//docs:razboj\_vsim\_raster\_test}.
\item \code{//soc:demo\_test} runs
\code{a\_core\_draws\_a\_solid\_and\_a\_gpu\_fills\_it}. Vreteno writes
a display list on the lattice, the rasteriser draws it, and every pixel
must match the model's picture of that list.
\end{itemize}

\subsection*{Used by}

\code{razboj::sim::run} and \code{razboj::sim::run\_list}, and so Razboj's
tests, \code{//gpu/razboj:trace} and \code{//gpu/razboj:demo}. The system in
\code{soc}, where the rasteriser is the host at corner 1,0 of the
lattice.

\subsection*{Limits}

\code{//docs:razboj} lists them. It writes one pixel per cycle at most. It
has no depth buffer, so primitives are drawn in list order. It has no
interpolation, so a triangle is one colour, and no texture. It clips
no geometry, only boxes, and the assembler \code{Op::encode} does that
on the host. It has one read out at a time.

The code adds these. It has no reset port, and phase 3 is final, so it
draws one list. \code{left} is eight bits. The pixel's offset from
\code{BASE} is formed at sixteen bits and then widened to \code{A}
bits. A colour is 24 bits.


\end{document}
